\subsection{RINGS OF FRACTIONS}

THEOREM 4. Let A be a commutative ring and S a subset of A. Let \(\mathbf{A}_{\mathrm{S}}\) be the monoid of fractions of A (provided only with multiplication) with denominators in S (§2, no. 4). Let \(\varepsilon: \mathbf{A} \to \mathbf{A}_{\mathrm{S}}\) be the canonical morphism. There exists on \(\mathbf{A}_{\mathrm{S}}\) one and only one addition satisfying the following conditions:

\begin{enumerate}
\def\labelenumi{(\alph{enumi})}
\item
  \(A_{s}\) , with this addition and its multiplication, is a commutative ring;
\item
  \(\varepsilon\) is a ring homomorphism.
\end{enumerate}

Suppose an addition has been found for \(\mathbf{A}_{\mathrm{s}}\) satisfying conditions (a) and (b).

Let \(x, y \in A_s\) . Let \(S'\) be the stable multiplicative submonoid of \(A\) generated by \(S\) . There exist \(a, b \in A\) and \(p, q \in S'\) such that \(x = a/p, y = b/q\) . Then

\[
x = \varepsilon (a q) \varepsilon (p q) ^ {- 1}, \quad y = \varepsilon (b p) \varepsilon (p q) ^ {- 1},
\]

whence

\begin{enumerate}
\def\labelenumi{(\arabic{enumi})}
\setcounter{enumi}{22}
\tightlist
\item
\end{enumerate}

\[
\begin{array}{r l} x + y & = (\varepsilon (a q) + \varepsilon (b p)) \varepsilon (p q) ^ {- 1} \\ & = \varepsilon (a q + b p) \varepsilon (p q) ^ {- 1} \\ & = (a q + b p) / p q. \end{array}
\]

This proves the uniqueness of the addition.

We now define an addition on \(A_{s}\) by setting \(x + y = (aq + bp)/pq\) . It is necessary to show that this definition does not depend on the choice of a, b, p, q. Now, if \(a', b' \in A\) , \(p', q' \in S'\) are such that \(x = a'/p', y = b'/q'\) , there exist s and t in \(S'\) such that \(ap's = a'ps\) , \(bq't = b'qt\) , whence

\[
(a q + b p) \left(p ^ {\prime} q ^ {\prime}\right) (s t) = \left(a ^ {\prime} q ^ {\prime} + b ^ {\prime} p ^ {\prime}\right) (p q) (s t)
\]

and hence

\[
(a q + b p) / p q = (a ^ {\prime} q ^ {\prime} + b ^ {\prime} p ^ {\prime}) / p ^ {\prime} q ^ {\prime}.
\]

It is easily verified that addition in \(A_{s}\) is associative and commutative, that 0/1 is identity element for addition, that \((-a)/p\) is the negative of a/p and that \(x(y + z) = xy + xz\) for all \(x, y, z \in A_{s}\) . If \(a, b \in A\) , then

\[
\varepsilon (a + b) = (a + b) / 1 = a / 1 + b / 1 = \varepsilon (a) + \varepsilon (b)
\]

and hence \(\varepsilon\) is a ring homomorphism.

DEFINITION 8. The ring defined in Theorem 4 is called the ring of fractions associated with S, or with denominators in S, and is denoted by A{[}S \(^{-1}\) {]}.

The zero of \(\mathbf{A}[\mathbf{S}^{-1}]\) is \(0 / 1\) , the unit of \(\mathbf{A}[\mathbf{S}^{-1}]\) is \(1 / 1\) .

We shall return to the properties of \(\mathbf{A}[\mathrm{S}^{-1}]\) in Commutative Algebra, Chapter II, § 2.

If S is the set of cancellable elements of A, the ring \(A[S^{-1}]\) is called the total ring of fractions of A. A is then identified with a subring of \(A[S^{-1}]\) by means of the mapping \(\varepsilon\) , which is then injective (I, §2, no. 4, Proposition 6).

THEOREM 5. Let A be a commutative ring, S a subset of A, B a ring and f a homomorphism of A into B such that every element of \(f(\mathbf{S})\) is invertible. There exists one and only one \(\bar{f}\) of \(A[S^{-1}]\) into B such that \(f = \bar{f} \circ \varepsilon\) .

We know (§ 2, no. 4, Theorem 1) that there exists one and only one morphism \(\bar{f}\) of the multiplicative monoid A{[}S \(^{-1}\) {]} into the multiplicative monoid B such that \(f = \bar{f} \circ \varepsilon\) . Let \(a, b \in \mathbf{A}, p, q \in \mathbf{S}'\) (stable multiplicative submonoid of A generated by S). As the elements of \(f(A)\) commute in pairs,

\[
\begin{array}{r l} \bar {f} (a | \dot {p} + b | q) & = \bar {f} ((a q + b \dot {p}) | \dot {p} q) = f (a q + b \dot {p}) f (\dot {p} q) ^ {- 1} \\ & = (f (a) f (q) + f (b) f (\dot {p})) f (\dot {p}) ^ {- 1} f (q) ^ {- 1} \\ & = f (a) f (\dot {p}) ^ {- 1} + f (b) f (q) ^ {- 1} \\ & = \bar {f} (a | \dot {p}) + \bar {f} (b | q). \end{array}
\]

Hence \(\tilde{f}\) is a ring homomorphism.

\section{FIELDS}

\subsection{FIELDS}

DEFINITION 1. A ring K is called a field if it does not consist only of 0 and every non-zero element of K is invertible.

The set of non-zero elements of the field K, with multiplication, is a group, which is precisely the multiplicative group \(K^{*}\) of the ring K ( \(\S\) 8, no. 1). The opposite ring of a field is a field. A field is called commutative if its multiplication is commutative; such a field is identified with its opposite. A non-commutative field is sometimes called a skew field.

Examples. *(1) We shall define in no. 4 the field of rational numbers; In General Topology the field of real numbers will be defined (General Topology, IV, § 1, no. 3), also the field of complex numbers (General Topology, VIII, § 1, no. 1) and the field of quaternions (General Topology, VIII, § 1, no. 4). These fields are commutative with the exception of the field of quaternions.* (2) The ring Z/2Z is obviously a field.

Let K be a field. Every subring L of K which is a field is called a subfield of K and K is then called an extension field of L; it amounts to the same to say that L is a subring of K and that \(x^{-1} \in L\) for every non-zero element x of L. If \((\mathbf{L}_{i})_{i \in I}\) is a family of subfields of K, then \(\bigcap_{i \in I} L_{i}\) is a subfield of K; for every subset X of K there thus exists a smallest subfield of K containing X; it is said to be generated by X.

PROPOSITION 1. Let K be a field. For every subset X of K, the centralizer (§ 8, no. 5, Example 3) X' of X is a subfield of K.

We know (loc. cit.) that \(\mathbf{X}'\) is a subring of \(\mathbf{K}\) . On the other hand, if \(x \neq 0\) is permutable with \(z \in \mathbf{X}\) , so is \(x^{-1}\) (I, §2, no. 3, Proposition 5) and hence \(\mathbf{X}'\) contains the inverse of every non-zero element of \(\mathbf{X}'\) .

COROLLARY. The centre of a field K is a (commutative) subfield of K.

THEOREM 1. Let A be a ring. The following conditions are equivalent:

\begin{enumerate}
\def\labelenumi{(\alph{enumi})}
\item
  A is a field;
\item
  A is not reduced to 0 and the only left ideals of A are 0 and A.
\end{enumerate}

Suppose that A is a field. Then A is not reduced to 0. Let a be a left ideal of A distinct from 0. There exists a non-zero a belonging to a. For all \(x \in A\) , \(x = (xa^{-1})a \in a\) ; hence a = A.

Suppose that A satisfies condition (b). Let \(x \neq 0\) be in A. We need to prove that x is invertible. The left ideal Ax contains x and hence is not zero, whence Ax = A. There thus exists \(x' \in A\) such that \(x'x = 1\) . Then \(x' \neq 0\) since \(1 \neq 0\) . Applying the above result to \(x'\) , it is seen that there exists \(x'' \in A\) such that \(x''x' = 1\) . Then \(x'' = x'' \cdot 1 = x''x'x = 1 \cdot x = x\) , hence \(xx' = 1\) . Hence \(x'\) is the inverse of x.

Remark. In Theorem 1, left ideals may be replaced by right ideals. In Chapter VIII, § 5, no. 2, we shall study non-zero rings A which have no two-sided ideal distinct from 0 and A; such rings (called quasi-simple) are not necessarily fields * (for example, the ring \(\mathbf{M}_2(\mathbf{Q})\) of square matrices of order 2 with rational coefficients is quasi-simple but is not a field)*.

COROLLARY 1. Let A be a ring and a two-sided ideal of A. For the ring A/a to be a field, it is necessary and sufficient that a be a maximal left ideal of A.

The left ideals of A/a are of the form b/a where b is a left ideal of A containing a (§ 8, no. 8, Corollary to Proposition 5). To say that A/a ≠ 0 means that a ≠ A. Under this hypothesis, A/a is a field if and only if the only left ideals of A containing a are a and A (Theorem 1), whence the corollary.

COROLLARY 2. Let A be a commutative ring which is not 0. There exists a homomorphism of A onto a commutative field.

By Krull's Theorem (§ 8, no. 6, Theorem 1), there exists in A a maximal ideal \(a\) . Then A/a is a field (Corollary 1).

COROLLARY 3. Let a be an integer \(\geqslant 0\) . For the ring \(\mathbf{Z} / a\mathbf{Z}\) to be a field, it is necessary and sufficient that a be prime.

This follows from Corollary 1 and § 8, no. 11.

For p prime the field Z/pZ is denoted by \(F_{p}\) .

THEOREM 2. Let K be a field and A a non-zero ring. If f is a homomorphism of K into A, then the subring \(f(\mathbf{K})\) of A is a field and f defines an isomorphism of K onto \(f(\mathbf{K})\) .

Let \(\mathfrak{a}\) be the kernel of \(f\) . Then \(1 \notin \mathfrak{a}\) for \(f(1) = 1 \neq 0\) in \(\mathbf{A}\) and, as \(\mathfrak{a}\) is a left ideal of \(\mathbf{K}\) , \(\mathfrak{a} = \{0\}\) by Theorem 1. Therefore \(f\) is injective and hence an isomorphism of \(\mathbf{K}\) onto the subring \(f(\mathbf{K})\) of \(\mathbf{A}\) ; the latter ring is therefore a field.

\subsection{INTEGRAL DOMAINS}

DEFINITION 2. A ring A is called an integral domain (or a domain of integrity) if it is commutative, non-zero, and the product of two non-zero elements of A is non-zero.

The ring Z of rational integers is an integral domain; it is commutative and non-zero; the product of two integers \textgreater0 is non-zero; every non-zero integer is of the form a or -a with a \textgreater{} 0 and \((-a)b = -ab\) , \((-a)(-b) = ab\) for a \textgreater{} 0, b \textgreater{} 0, whence our assertion.

Every commutative field is an integral domain. A subring of an integral domain is an integral domain. In particular, a subring of a commutative field is an integral domain. We shall show that conversely every integral domain A is isomorphic to a subring of a commutative field. Recall ( \(\S 8\) , no. 12) that A has been identified with a subring of its total ring of fractions. Our assertion then follows from the following proposition:

PROPOSITION 2. If A is an integral domain, the total ring of fractions K of A is a commutative ring.

The ring K is commutative. It is non-zero since \(A \neq \{0\}\) . As A is an integral domain, every non-zero element of A is cancellable and K consists of the fractions a/b with \(b \neq 0\) . Now \(a/b \neq 0\) implies \(a \neq 0\) and the fraction b/a is then inverse of a/b.

The total ring of fractions of an integral domain is called its field of fractions. Such a ring is identified with its image in its field of fractions.

PROPOSITION 3. Let B be a non-zero ring and A a commutative subring of B such that every non-zero element of A is invertible in B.

\begin{enumerate}
\def\labelenumi{(\alph{enumi})}
\item
  A is an integral domain.
\item
  Let \(\mathbf{A}'\) be the field of fractions \(\mathbf{A}\) . The canonical injection of \(\mathbf{A}\) into \(\mathbf{B}\) can be extended uniquely to an isomorphism \(f\) of \(\mathbf{A}'\) onto a subfield of \(\mathbf{B}\) .
\item
  The elements of \(f(\mathbf{A}')\) are the \(xy^{-1}\) where \(x \in \mathbf{A}, y \in \mathbf{A}, y \neq 0\) .
\end{enumerate}

Assertion (a) is obvious. The canonical injection of A into B extends uniquely to a homomorphism \(f\) of \(\mathbf{A}'\) into \(\mathbf{B}\) (§ 8, no. 12, Theorem 5). Assertion (b) then follows from no. 1, Theorem 2. The elements of \(\mathbf{A}'\) are the fractions \(x / y\) with \(x \in \mathbf{A}, y \in \mathbf{A}, y \neq 0\) and \(f(x / y) = xy^{-1}\) , whence (c).

\subsection{PRIME IDEALS}

PROPOSITION 4. Let A be a commutative ring and p an ideal of A. The following conditions are equivalent:

\begin{enumerate}
\def\labelenumi{(\alph{enumi})}
\item
  the ring A/p is an integral domain;
\item
  \(\mathbf{A} \neq \mathfrak{p}\) and, if \(x \in \mathbf{A} - \mathfrak{p}\) and \(y \in \mathbf{A} - \mathfrak{p}\) , then \(xy \in \mathbf{A} - \mathfrak{p}\) .
\item
  \(\mathfrak{p}\) is the kernel of a homomorphism of \(\mathbf{A}\) into a field.
\end{enumerate}

The implications (c) \(\Rightarrow\) (b) \(\Rightarrow\) (a) are obvious. If A/p is an integral domain, let \(f\) be the canonical injection of A/p into its field of fractions and \(g\) the canonical homomorphism of A onto A/p; then p is the kernel of \(f \circ g\) , whence the implication (a) \(\Rightarrow\) (c).

DEFINITION 3. In a commutative ring A, an ideal p satisfying the conditions of Proposition 4 is called a prime ideal.

Examples. (1) Let A be a commutative ring. If m is a maximal ideal of A, m is prime; the ring A/m is a field (no. 1, Corollary 1 to Theorem 1).

\begin{enumerate}
\def\labelenumi{(\arabic{enumi})}
\setcounter{enumi}{1}
\tightlist
\item
  If \(\mathbf{A}\) is an integral domain, the ideal \(\{0\}\) of \(\mathbf{A}\) is prime (but not maximal in general, as the example of the ring \(\mathbf{Z}\) proves).
\end{enumerate}

\subsection{THE FIELD OF RATIONAL NUMBERS}

DEFINITION 4. The field of fractions of the ring Z of rational integers is called the field of rational numbers and is denoted by Q. The elements of Q are called rational numbers.

Every rational number is thus of the form \(a / b\) where \(a\) and \(b\) are rational integers with \(b \neq 0\) (and we may even take \(b > 0\) as the relation

\[
a / b = (- a) / (- b)
\]

proves). \(\mathbf{Q}_{+}\) is used to denote the set of rational numbers of the form \(a / b\) with \(a\in \mathbf{N}\) and \(b\in \mathbf{N}^*\) .

We have the relations:

\begin{enumerate}
\def\labelenumi{(\arabic{enumi})}
\tightlist
\item
\end{enumerate}

\[
\mathbf {Q} _ {+} + \mathbf {Q} _ {+} = \mathbf {Q} _ {+}\tag{2}
\]

\[
\mathbf {Q} _ {+}. \mathbf {Q} _ {+} = \mathbf {Q} _ {+}\tag{3}
\]

\[
\mathbf {Q} _ {+} \cap (- \mathbf {Q} _ {+}) = \{0 \}\tag{4}
\]

\[
\mathbf {Q} _ {+} \cup (- \mathbf {Q} _ {+}) = \mathbf {Q}\tag{5}
\]

\[
\mathbf {Q} _ {+} \cap \mathbf {Z} = \mathbf {N}.
\]

The first two follow from the formulae \(a / b + a' / b' = (ab' + a'b) / bb'\) , \((a / b)(a'/b') = aa'/bb'\) , \(0 \in \mathbf{Q}_+\) , \(1 \in \mathbf{Q}_+\) and the fact that \(\mathbf{N}\) is stable under addition and multiplication and \(\mathbf{N}^*\) under multiplication. Obviously \(0 \in \mathbf{Q}_+\) , whence \(0 \in (-\mathbf{Q}_+)\) ; let \(x\) be in \(\mathbf{Q}_+ \cap (-\mathbf{Q}_+)\) ; then there exist positive integers \(a, b, a', b'\) with \(b \neq 0\) , \(b' \neq 0\) and \(x = a / b = -a'/b'\) ; then \(ab' + ba' = 0\) , whence \(ab' = 0\) (for \(ab' \geqslant 0\) and \(ba' \geqslant 0\) ) and therefore \(a = 0\) since \(b' \neq 0\) ; in other words, \(x = 0\) . Finally, obviously \(\mathbf{N} \subset \mathbf{Z}\) and \(\mathbf{N} \subset \mathbf{Q}_+\) . Conversely, if \(x\) belongs to \(\mathbf{Z} \cap \mathbf{Q}_+\) , it is a rational integer; there exist two rational integers \(a\) and \(b\) with \(a \geqslant 0\) , \(b > 0\) and \(x = a / b\) , whence \(a = bx\) ; if \(x \notin \mathbf{N}\) , then \(-x > 0\) , whence \(-a = b(-x) > 0\) and consequently \(a < 0\) contrary to the hypothesis.

Given two rational numbers \(x\) and \(y\) , we write \(x \leqslant y\) if \(y - x \in \mathbf{Q}_+\) . It is easily deduced from formulae (1), (3) and (4) that \(x \leqslant y\) is a total ordering on

Q, from (5) that this relation induces the usual order relation on Z. Finally, from (1) it follows that the relations \(x \leqslant y\) and \(x' \leqslant y'\) imply \(x + x' \leqslant y + y'\) and from (2) that the relation \(x \leqslant y\) implies \(xz \leqslant yz\) for all \(z \geqslant 0\) and \(xz \geqslant yz\) for \(z \leqslant 0\) (cf.~VI, § 2, no. 1).

Let \(x\) be a rational number. \(x\) is said to be positive if \(x \geqslant 0\) , strictly positive if \(x > 0\) , negative if \(x \leqslant 0\) and strictly negative if \(x < 0\) . The set of positive rational numbers is \(\mathbf{Q}_{+}\) and that of negative numbers is \(-\mathbf{Q}_{+}\) . If \(\mathbf{Q}^{*}\) denotes the set of non-zero rational numbers, the set \(\mathbf{Q}_{+}^{*}\) of strictly positive numbers is equal to \(\mathbf{Q}^{*} \cap \mathbf{Q}_{+}\) and \(-\mathbf{Q}_{+}^{*}\) is the set of strictly negative numbers.

The sets \(\mathbf{Q}_{+}^{*}\) and \(\{1, -1\}\) are subgroups of the multiplicative group \(\mathbf{Q}^{*}\) . Every rational number \(x \neq 0\) can be expressed in one and only one way in the form \(1.y\) , \((-1).y\) , where \(y > 0\) ; hence the multiplicative group \(\mathbf{Q}^{*}\) is the product of the subgroups \(\mathbf{Q}_{+}^{*}\) and \(\{1, -1\}\) , the component of \(x\) in \(\mathbf{Q}_{+}^{*}\) is called the absolute value of \(x\) and is denoted by \(|x|\) ; the component of \(x\) in \(\{-1, 1\}\) (equal to 1 if \(x > 0\) , to \(-1\) if \(x < 0\) ) is called the sign of \(x\) and denoted by \(\operatorname{sgn} x\) .

Usually these two functions are extended to the whole of \(\mathbf{Q}\) by setting \(|0| = 0\) and \(\operatorname{sgn} 0 = 0\) .

\section{INVERSE AND DIRECT LIMITS}

Throughout this paragraph, I will denote a non-empty preordered set, \(\alpha \leqslant \beta\) the preordering on I. The notion of inverse (resp. direct) system of sets relative to the indexing set I is defined in Set Theory, III, § 7, no. 1 (resp. Set Theory, III, § 7, no. 5, under the hypothesis that I is right directed).

\subsection{INVERSE SYSTEMS OF MAGMAS}

DEFINITION 1. An inverse system of magmas relative to the indexing set I is an inverse system of sets \((\mathbf{E}_{\alpha}, f_{\alpha\beta})\) relative to I, each \(E_{\alpha}\) having a magma structure and each \(f_{\alpha\beta}\) being a magma homomorphism.

Let \((\mathbf{E}_{\alpha}, f_{\alpha \beta})\) be an inverse system of magmas whose laws are written multiplicatively. The inverse limit set \(\mathbf{E} = \varprojlim \mathbf{E}_{\alpha}\) is a subset of the product magma \(\prod_{\alpha \in I} \mathbf{E}_{\alpha}\) consisting of the families \((x_{\alpha})_{\alpha \in I}\) such that \(x_{\alpha} = f_{\alpha \beta}(x_{\beta})\) for \(\alpha \leqslant \beta\) . If \((x_{\alpha})\) and \((y_{\alpha})\) belong to E, then for \(\alpha \leqslant \beta\) , \(x_{\alpha} = f_{\alpha \beta}(x_{\beta})\) and \(y_{\alpha} = f_{\alpha \beta}(y_{\beta})\) , hence \(x_{\alpha}y_{\alpha} = f_{\alpha \beta}(x_{\beta})f_{\alpha \beta}(y_{\beta}) = f_{\alpha \beta}(x_{\beta}y_{\beta})\) ; hence E is a submagma of \(\prod_{\alpha \in I} \mathbf{E}_{\alpha}\) . E will be given the law induced by that on \(\prod_{\alpha \in I} \mathbf{E}_{\alpha}\) ; the magma obtained is called the inverse limit magma of the magmas \(E_{\alpha}\) . It enjoys the following universal property:

\begin{enumerate}
\def\labelenumi{(\alph{enumi})}
\item
  For all \(\alpha \in \mathbf{I}\) , the canonical mapping \(f_{\alpha}\) of E into \(\mathrm{E}_{\alpha}\) is a magma homomorphism of E into \(\mathrm{E}_{\alpha} \cdot f_{\alpha} = f_{\alpha \beta} \circ f_{\beta}\) for \(\alpha \leqslant \beta\) .
\item
  Suppose a magma F is given and homomorphisms \(u_{\alpha} \colon \mathrm{F} \to \mathrm{E}_{\alpha}\) such that \(u_{\alpha} = f_{\alpha \beta} \circ u_{\beta}\) for \(\alpha \leqslant \beta\) . There exists one and only one homomorphism \(u \colon \mathrm{F} \to \mathrm{E}\) such that \(u_{\alpha} = f_{\alpha} \circ u\) for all \(\alpha \in \mathrm{I}\) (namely \(x \mapsto u(x) = (u_{\alpha}(x))_{\alpha \in \mathrm{I}}\) ).
\end{enumerate}

If the magmas \(E_{\alpha}\) are associative (resp. commutative), so is E. Suppose that each magma \(E_{\alpha}\) admits an identity element \(e_{\alpha}\) and that the homomorphisms \(f_{\alpha\beta}\) are unital. Then \(e = (e_{\alpha})_{\alpha \in I}\) belongs to E for \(e_{\alpha} = f_{\alpha\beta}(e_{\beta})\) for \(\alpha \leqslant \beta\) and it is an identity element of the magma E; with the above notation, the homomorphisms \(f_{\alpha}\) are unital and if the \(u_{\alpha}\) are unital then u is unital. Further, an element \(x = (x_{\alpha})_{\alpha \in I}\) of E is invertible if and only if each of the \(x_{\alpha}\) is invertible in the corresponding magma \(E_{\alpha}\) and \(x^{-1} = (x_{\alpha}^{-1})_{\alpha \in I}\) ; this follows from the formula \(f_{\alpha\beta}(x_{\beta}^{-1}) = f_{\alpha\beta}(x_{\beta})^{-1} = x_{\alpha}^{-1}\) for \(\alpha \leqslant \beta\) .

From these remarks it can be deduced that if the magmas \(E_{\alpha}\) are monoids (resp. groups) and the \(f_{\alpha\beta}\) are monoid homomorphisms, then the magma E is a monoid (resp. a group). In this case we speak of an inverse system of monoids (resp. groups). The universal property goes over immediately to this case.

It is left to the reader to define an inverse system of rings \((\mathrm{E}_{\alpha}, f_{\alpha \beta})\) and to verify that \(\mathbf{E} = \varprojlim \mathbf{E}_{\alpha}\) is a subring of the product ring \(\prod_{\alpha \in I} \mathbf{E}_{\alpha}\) , called the inverse limit ring of the rings \(\mathbf{E}_{\alpha}\) ; it can be verified that the universal property extends to this case.

Let \(\mathfrak{E} = (\mathrm{E}_{\alpha}, f_{\alpha\beta})\) and \(\mathfrak{E}' = (\mathrm{E}_{\alpha}', f_{\alpha\beta}')\) be two inverse systems of magmas (resp. monoids, groups, rings) relative to the same indexing set. A homomorphism of \(\mathfrak{E}\) into \(\mathfrak{E}'\) is an inverse system \((u_{\alpha})_{\alpha \in I}\) of mappings \(u_{\alpha}: \mathrm{E}_{\alpha} \to \mathrm{E}_{\alpha}'\) such that each \(u_{\alpha}\) is a homomorphism. Under these conditions, the mapping \(u = \varprojlim u_{\alpha}\) of \(\varprojlim \mathrm{E}_{\alpha}\) into \(\varprojlim \mathrm{E}_{\alpha}'\) is a homomorphism (cf.~Set Theory, III, § 7, no. 2).

\subsection{INVERSE LIMITS OF ACTIONS}

Suppose there are given two inverse systems of sets \((\Omega_{\alpha}, \phi_{\alpha\beta})\) and \((\mathrm{E}_{\alpha}, f_{\alpha\beta})\) relative to the same indexing set I. Suppose there is given for all \(\alpha \in I\) an action of \(\Omega_{\alpha}\) on \(E_{\alpha}\) such that

\[
f _ {\alpha \beta} (\omega_ {\beta} x _ {\beta}) = \phi_ {\alpha \beta} (\omega_ {\beta}) \cdot f _ {\alpha \beta} (x _ {\beta})\tag{1}
\]

for \(\alpha \leqslant \beta, x_{\beta} \in \mathrm{E}_{\beta}, \omega_{\beta} \in \Omega_{\beta}\) . Then the family of actions considered is said to be an inverse system of actions. Let \(\Omega = \varprojlim \Omega_{\alpha}\) and \(\mathrm{E} = \varprojlim \mathrm{E}_{\alpha}\) ; if \(x = (x_{\alpha})_{\alpha \in I}\) belongs to \(\mathrm{E}\) and \(\omega = (\omega_{\alpha})_{\alpha \in I}\) belongs to \(\Omega\) , then \(\omega \cdot x = (\omega_{\alpha} \cdot x_{\alpha})_{\alpha \in I}\) belongs to E by (1). Thus an action of \(\Omega\) on E is defined called the inverse limit of the actions of the \(\Omega_{\alpha}\) on the \(\mathrm{E}_{\alpha}\) .

The above applies especially in the case where the \(\Omega_{\alpha}\) are monoids and each action of \(\Omega_{\alpha}\) on \(E_{\alpha}\) is an operation. Then the inverse limit of these operations is an operation of the monoid on E.

It is left to the reader to define the inverse limit of an inverse system of groups with operators and to verify that this limit is a group with operators.

\subsection{DIRECT SYSTEMS OF MAGMAS}

In this no. and the following we shall assume that I is right directed.

DEFINITION 2. A direct system of magmas relative to the indexing set I is a direct system of sets \((\mathrm{E}_{\alpha}, f_{\beta\alpha})\) relative to I, each \(E_{\alpha}\) having a magma structure and each \(f_{\beta\alpha}\) being a magma homomorphism.

Let \((\mathrm{E}_{\alpha}, f_{\beta\alpha})\) be a direct system of magmas. E will denote the direct limit set \(\lim_{\rightarrow} \mathrm{E}_{\alpha}\) and \(f_{\alpha}\) the canonical mapping of \(\mathrm{E}_{\alpha}\) into E. Recall that

\begin{enumerate}
\def\labelenumi{(\arabic{enumi})}
\setcounter{enumi}{1}
\tightlist
\item
\end{enumerate}

\[
f _ {\beta} \circ f _ {\beta \alpha} = f _ {\alpha} \quad \text { for } \alpha \leqslant \beta ,\tag{3}
\]

\[
\mathrm{E} = \bigcup_ {\alpha \in \mathbf {I}} f _ {\alpha} (\mathrm{E} _ {\alpha}).
\]

By (2), also

\[
f _ {\alpha} (\mathrm{E} _ {\alpha}) \subset f _ {\beta} (\mathrm{E} _ {\beta}) \quad \text { for } \alpha \leqslant \beta .\tag{4}
\]

If \(x_{\alpha}, y_{\alpha} \in \mathrm{E}_{\alpha}\) are such that \(f_{\alpha}(x_{\alpha}) = f_{\alpha}(y_{\alpha})\) , there exists a \(\beta \geqslant \alpha\) such that \(f_{\beta \alpha}(x_{\alpha}) = f_{\beta \alpha}(y_{\alpha})\) .

PROPOSITION 1. There exists on E one and only one magma structure for which the mappings \(f_{\alpha} \colon \mathrm{E}_{\alpha} \to \mathrm{E}\) are homomorphisms. If the magmas \(\mathrm{E}_{\alpha}\) are associative (resp. commutative), so is E. If the magmas \(\mathrm{E}_{\alpha}\) and the homomorphism \(f_{\beta \alpha}\) are unital, so are the magma E and the homomorphisms \(f_{\alpha}\) .

The magmas \(E_{\alpha}\) will be written multiplicatively.

Let \(x, y\) be in E. There exist \(\alpha\) in I and \(x_{\alpha}, y_{\alpha}\) in \(E_{\alpha}\) such that \(x = f_{\alpha}(x_{\alpha})\) and \(y = f_{\alpha}(y_{\alpha})\) . If there exists a magma structure on E for which \(f_{\alpha}\) is a homomorphism, then \(x.y = f_{\alpha}(x_{\alpha}y_{\alpha})\) , whence the uniqueness of this magma structure.

To prove the existence, we must prove that for \(\alpha, \beta\) in I, \(x_{\alpha}, y_{\alpha}\) in E \(_{\alpha}\) and \(x_{\beta}', y_{\beta}'\) in E \(_{\beta}\) , the relations

\[
f _ {\alpha} (x _ {\alpha}) = f _ {\beta} (x _ {\beta} ^ {\prime}), \quad f _ {\alpha} (y _ {\alpha}) = f _ {\beta} (y _ {\beta} ^ {\prime})\tag{5}
\]

imply \(f_{\alpha}(x_{\alpha}y_{\alpha}) = f_{\beta}(x_{\beta}'y_{\beta}')\) . For \(\gamma \geqslant \alpha\) and \(\gamma \geqslant \beta\) , we set \(x_{\gamma} = f_{\gamma \alpha}(x_{\alpha})\) , \(y_{\gamma} = f_{\gamma \alpha}(y_{\alpha}), x_{\gamma}^{\prime} = f_{\gamma \beta}(x_{\beta}^{\prime}), y_{\gamma}^{\prime} = f_{\gamma \beta}(y_{\beta}^{\prime})\) . By the definition of direct limit, there exists \(\gamma\) in I with \(\gamma \geqslant \alpha, \gamma \geqslant \beta, x_{\gamma} = x_{\gamma}^{\prime}, y_{\gamma} = y_{\gamma}^{\prime}\) . Then

\[
\begin{array}{r l} f _ {\alpha} (x _ {\alpha} y _ {\alpha}) & = f _ {\gamma} (f _ {\gamma \alpha} (x _ {\alpha} y _ {\alpha})) = f _ {\gamma} (x _ {\gamma} y _ {\gamma}) = f _ {\gamma} (x _ {\gamma} ^ {\prime} y _ {\gamma} ^ {\prime}) = f _ {\gamma} (f _ {\gamma \beta} (x _ {\beta} ^ {\prime} y _ {\beta} ^ {\prime})) \\ & = f _ {\beta} (x _ {\beta} ^ {\prime} y _ {\beta} ^ {\prime}). \end{array}
\]

Suppose the magmas \(E_{\alpha}\) are associative. Let x, y, z be in E. There exist \(\alpha \in I\) and elements \(x_{\alpha}, y_{\alpha}, z_{\alpha}\) in \(E_{\alpha}\) such that

\[
x = f _ {\alpha} (x _ {\alpha}), \quad y = f _ {\alpha} (y _ {\alpha}), \quad z = f _ {\alpha} (z _ {\alpha}).
\]

Then \(xy = f_{\alpha}(x_{\alpha}y_{\alpha})\) , whence \((xy)z = f_{\alpha}((x_{\alpha}y_{\alpha})z_{\alpha})\) ; similarly

\[
x (y z) = f _ {\alpha} \left(x _ {\alpha} \left(y _ {\alpha} z _ {\alpha}\right)\right),
\]

whence \((xy)z = x(yz)\) for \((x_{\alpha}y_{\alpha})z_{\alpha} = x_{\alpha}(y_{\alpha}z_{\alpha})\) . The case of commutative magmas is treated analogously.

Suppose finally that each magma \(E_{\alpha}\) has an identity element \(e_{\alpha}\) and that \(f_{\beta\alpha}(e_{\alpha}) = e_{\beta}\) for \(\alpha \leqslant \beta\) . For \(\alpha, \beta\) in I, there exists \(\gamma\) in I with \(\gamma \geqslant \alpha\) and \(\gamma \geqslant \beta\) , whence

\[
f _ {\alpha} (e _ {\alpha}) = f _ {\gamma} \left(f _ {\gamma \alpha} (e _ {\alpha})\right) = f _ {\gamma} (e _ {\gamma}) = f _ {\gamma} \left(f _ {\gamma \beta} (e _ {\beta})\right) = f _ {\beta} (e _ {\beta})
\]

and there thus exists an element \(e\) in \(\mathbf{E}\) such that \(f_{\alpha}(e_{\alpha}) = e\) for all \(\alpha \in \mathbf{I}\) . Let \(x\) be in \(\mathbf{E}\) ; choose \(\alpha \in \mathbf{I}\) and \(x_{\alpha} \in \mathrm{E}_{\alpha}\) such that \(x = f_{\alpha}(x_{\alpha})\) . Then

\[
e x = f _ {\alpha} (e _ {\alpha}) \cdot f _ {\alpha} (x _ {\alpha}) = f _ {\alpha} (e _ {\alpha}. x _ {\alpha}) = f _ {\alpha} (x _ {\alpha}) = x
\]

and similarly \(x.e = x\) , hence \(e\) is the identity element of \(E\) .

The magma E is called the direct limit of the magmas \(E_{\alpha}\) .

PROPOSITION 2. Let \((\mathrm{E}_{\alpha}, f_{\beta \alpha})\) be a direct system of magmas and let E be its direct limits \(f_{\alpha} \colon \mathrm{E}_{\alpha} \to \mathrm{E}\) the canonical homomorphisms. Suppose a magma F and homomorphisms \(u_{\alpha} \colon \mathrm{E}_{\alpha} \to \mathrm{F}\) are given such that \(u_{\alpha} = u_{\beta} \circ f_{\beta \alpha}\) for \(\alpha \leqslant \beta\) . There exists one and only one homomorphism \(u \colon \mathrm{E} \to \mathrm{F}\) such that \(u_{\alpha} = u \circ f_{\alpha}\) for all \(\alpha \in \mathbf{I}\) . If the magmas \(\mathrm{E}_{\alpha}\) and F and the homomorphisms \(f_{\beta \alpha}\) and \(u_{\alpha}\) are unital, the homomorphism \(u\) is unital.

We know (Set Theory, III, § 7, no. 6, Proposition 6) that there exists one and only one mapping \(u: \mathrm{E} \to \mathrm{F}\) such that \(u_{\alpha} = u \circ f_{\alpha}\) for all \(\alpha \in \mathbf{I}\) . We verify that \(u\) is a homomorphism: let \(x, y\) be in \(\mathrm{E}, \alpha\) in \(\mathrm{I}\) and \(x_{\alpha}, y_{\alpha}\) in \(\mathrm{E}_{\alpha}\) such that \(x = f_{\alpha}(x_{\alpha})\) and \(y = f_{\alpha}(y_{\alpha})\) . Then \(xy = f_{\alpha}(x_{\alpha}y_{\alpha})\) , whence

\[
\begin{array}{r l} u (x y) & = u (f _ {\alpha} (x _ {\alpha} y _ {\alpha})) = u _ {\alpha} (x _ {\alpha} y _ {\alpha}) = u _ {\alpha} (x _ {\alpha}) u _ {\alpha} (y _ {\alpha}) \\ & = u (f _ {\alpha} (u _ {\alpha})) u (f _ {\alpha} (y _ {\alpha})) = u (x) u (y). \end{array}
\]

We consider now the unital case and let \(e_{\alpha}\) denote the identity element of \(E_{\alpha}\) , \(e\) that of \(E\) and \(e'\) that of \(F\) . Choose \(\alpha \in I\) , then \(e = f_{\alpha}(e_{\alpha})\) , whence

\[
u (e) = u \left(f _ {\alpha} \left(e _ {\alpha}\right)\right) = u _ {\alpha} \left(e _ {\alpha}\right) = e ^ {\prime}
\]

for \(u_{\alpha}\) is unital. Hence \(u\) is unital.

By analogy with the notion of direct system of magmas, that of a direct system of monoids or groups can be formulated. Proposition 1 shows that the magma E which is the limit of a direct system of monoids \((\mathrm{E}_{\alpha}, f_{\beta \alpha})_{\alpha, \beta \in I}\) is a monoid. We show that E is a group if the \(\mathrm{E}_{\alpha}\) are groups: let \(x \in \mathrm{E}\) , \(\alpha \in I\) and \(x_{\alpha} \in \mathrm{E}_{\alpha}\) be such that \(x = f_{\alpha}(x_{\alpha})\) ; the element \(y = f_{\alpha}(x_{\alpha}^{-1})\) of E is the inverse of \(x\) (§ 2, no. 3). The universal property of Proposition 2 goes over immediately in the case of a direct system of monoids or groups.

The reader is left to define a direct system of rings. Let \((\mathbf{A}_{\alpha}, f_{\beta \alpha})\) be such a direct system; let \(\mathbf{A} = \lim_{\longrightarrow} \mathbf{A}_{\alpha}\) and \(f_{\alpha}: \mathbf{A}_{\alpha} \to \mathbf{A}\) the canonical homomorphisms. There exists (Proposition 2) on \(\mathbf{A}\) one and only one addition and multiplication characterized by \(x + y = f_{\alpha}(x_{\alpha} + y_{\alpha})\) , \(xy = f_{\alpha}(x_{\alpha}y_{\alpha})\) for \(\alpha \in \mathbf{I}\) , \(x_{\alpha}, y_{\alpha}\) in \(\mathbf{A}_{\alpha}\) and \(x = f_{\alpha}(x_{\alpha})\) , \(y = f_{\alpha}(y_{\alpha})\) . Under addition \(\mathbf{A}\) is a commutative group and multiplication is associative and unital. Finally, for \(x, y, z\) in \(\mathbf{A}\) , choose \(\alpha\) in \(\mathbf{I}\) and \(x_{\alpha}, y_{\alpha}\) and \(z_{\alpha}\) in \(\mathbf{A}_{\alpha}\) with

\[
x = f _ {\alpha} (x _ {\alpha}), \quad y = f _ {\alpha} (y _ {\alpha}) \quad \text { and } \quad z = f _ {\alpha} (z _ {\alpha}).
\]

Then

\[
\begin{array}{r l} (x + y) \cdot z & = f _ {\alpha} (x _ {\alpha} + y _ {\alpha}) f _ {\alpha} (z _ {\alpha}) = f _ {\alpha} ((x _ {\alpha} + y _ {\alpha}) z _ {\alpha}) \\ & = f _ {\alpha} (x _ {\alpha} z _ {\alpha} + y _ {\alpha} z _ {\alpha}) = f _ {\alpha} (x _ {\alpha} z _ {\alpha}) + f _ {\alpha} (y _ {\alpha} z _ {\alpha}) = x z + y z \end{array}
\]

and the relation \(x(y + z) = xy + xz\) is proved analogously. In other words, A has the structure of a ring, characterized by the fact that \(f_{\alpha}\) is a ring homomorphism for all \(\alpha \in I\) .

The ring A is called the direct limit of the rings \(A_{\alpha}\) . Proposition 2 extends immediately to the case of rings.

PROPOSITION 3. (a) If the \(\mathbf{A}_{\alpha}\) are non-zero, \(\mathbf{A}\) is non-zero.

\begin{enumerate}
\def\labelenumi{(\alph{enumi})}
\setcounter{enumi}{1}
\item
  If the \(\mathbf{A}_{\alpha}\) are integral domains, \(\mathbf{A}\) is an integral domain.
\item
  If the \(\mathbf{A}_{\alpha}\) are fields, \(\mathbf{A}\) is a field.
\end{enumerate}

Let \(0_{\alpha}, 1_{\alpha}\) be the zero and unit of \(A_{\alpha}\) and 0, 1 the zero and unit of A. There exists \(\alpha \in I\) such that \(f_{\alpha}(0_{\alpha}) = 0, f_{\alpha}(1_{\alpha}) = 1\) . If 0 = 1, there exists \(\beta \geqslant \alpha\) such that \(f_{\beta \alpha}(0_{\alpha}) = f_{\beta \alpha}(1_{\alpha})\) , that is \(0_{\beta} = 1_{\beta}\) . This proves (a).

Suppose that \(A_{\alpha}\) are integral domains. Then A is commutative and non-zero by (a). Let x, y be elements of A such that xy = 0. There exists \(\alpha \in I\) and \(x_{\alpha}, y_{\alpha} \in A_{\alpha}\) such that \(x = f_{\alpha}(x_{\alpha}), y = f_{\alpha}(y_{\alpha})\) . Then \(f_{\alpha}(x_{\alpha}y_{\alpha}) = xy = 0 = f_{\alpha}(0_{\alpha})\) . Hence there exists \(\beta \geqslant \alpha\) such that \(f_{\beta\alpha}(x_{\alpha}y_{\alpha}) = f_{\beta\alpha}(0_{\alpha})\) . As \(A_{\beta}\) is an integral domain, it follows that \(f_{\beta\alpha}(x_{\alpha}) = 0_{\beta}\) or \(f_{\beta\alpha}(y_{\alpha}) = 0_{\beta}\) , hence x = 0 or y = 0. This proves (b).

Suppose that the \(A_{\alpha}\) are fields. Then \(A \neq \{0\}\) by (a). Let x be a non-zero element of A. There exist \(\alpha \in I\) and \(x_{\alpha} \in A_{\alpha}\) such that \(x = f_{\alpha}(x_{\alpha})\) . Then \(x_{\alpha} \neq 0\) and \(f_{\alpha}(x_{\alpha}^{-1})\) is the inverse of x in A. This proves (c).

Let \(\mathfrak{E} = (\mathrm{E}_{\alpha}, f_{\beta \alpha})\) and \(\mathfrak{E}' = (\mathrm{E}_{\alpha}', f_{\beta \alpha}')\) be two direct systems of magmas (resp. monoids, groups, rings). A homomorphism of \(\mathfrak{E}\) into \(\mathfrak{E}'\) is a direct system \((u_{\alpha})_{\alpha \in I}\) of mappings \(u_{\alpha} \colon E_{\alpha} \to E_{\alpha}'\) such that each \(u_{\alpha}\) is a homomorphism. Under these conditions, the mapping \(u = \lim u_{\alpha}\) from \(E = \lim E_{\alpha}\) to \(E' = \lim E_{\alpha}'\) is a homomorphism (cf.~Set Theory, III, § 7, no. 6).

\subsection{DIRECT LIMIT OF ACTIONS}

Suppose that there are given two direct systems of sets \((\Omega_{\alpha}, \phi_{\beta\alpha})\) and \((\mathrm{E}_{\alpha}, f_{\beta\alpha})\) relative to the same indexing set I and for each \(\alpha \in I\) an action of \(\Omega_{\alpha}\) on \(E_{\alpha}\) . Suppose that

\[
f _ {\beta \alpha} (\omega_ {\alpha}. x _ {\alpha}) = \phi_ {\beta \alpha} (\omega_ {\alpha}). f _ {\beta \alpha} (x _ {\alpha})
\]

for \(\alpha \leqslant \beta\) , \(\omega_{\alpha} \in \Omega_{\alpha}\) and \(x_{\alpha} \in \mathrm{E}_{\alpha}\) . Then the family of actions under consideration is said to be a direct system of actions. It is easily verified as in Proposition 2 that there exists an action \(h\) of \(\Omega = \varinjlim \Omega_{\alpha}\) on \(\mathrm{E} = \varinjlim \mathrm{E}_{\alpha}\) which is described as follows: let \(\omega \in \Omega\) and \(x \in \mathrm{E}\) ; choose \(\alpha \in \mathrm{I}\) and \(\omega_{\alpha} \in \Omega_{\alpha}\) , \(x_{\alpha} \in \mathrm{E}_{\alpha}\) such that \(\omega = \phi_{\alpha}(\omega_{\alpha})\) and \(x = f_{\alpha}(x_{\alpha})\) ( \(\phi_{\alpha}: \Omega_{\alpha} \to \Omega\) and \(f_{\alpha}: \mathrm{E}_{\alpha} \to \mathrm{E}\) denote the canonical mappings); then \(\omega . x = f_{\alpha}(\omega_{\alpha}. x_{\alpha})\) . The action of \(\Omega\) on \(\mathrm{E}\) is called the direct limit of the actions of the \(\Omega_{\alpha}\) on the \(\mathrm{E}_{\alpha}\) .

If the \(\Omega_{\alpha}\) are monoids and each action of \(\Omega_{\alpha}\) on \(E_{\alpha}\) is an operation, the direct limit action is an operation.

It is left to the reader to define the direct limit of a direct system of groups with operators and to verify that this limit is a group with operators.

\BourbakiExercisesSection

\begin{enumerate}
\def\labelenumi{\arabic{enumi}.}
\tightlist
\item
  Let E be a set, \(A \subset E \times E\) and \((x, y) \mapsto x \top y\) , \((x, y) \in A\) , an internal law of composition not everywhere defined on E. Given two subsets X and Y of E, let \(X \top Y\) denote the set of elements of E of the form \(x \top y\) with \(x \in X\) , \(y \in Y\) and \((x, y) \in A\) . A law of composition is thus defined which is everywhere defined on \(\mathfrak{P}(E)\) .
\end{enumerate}

If \((\mathbf{X}_{\alpha})_{\alpha \in \mathbf{A}}\) and \((\mathbf{Y}_{\beta})_{\beta \in B}\) are any two families of subsets of E, then

\[
\left(\bigcup_ {\alpha \in A} X _ {\alpha}\right) \top \left(\bigcup_ {\beta \in B} Y _ {\beta}\right) = \bigcup_ {(\alpha , \beta) \in A \times B} (X _ {\alpha} \top Y _ {\beta}).
\]

\begin{enumerate}
\def\labelenumi{\arabic{enumi}.}
\setcounter{enumi}{1}
\item
  Let \(\top\) be a law of composition not everywhere defined on a set E. Let E' be the subset of \(\mathfrak{P}(E)\) consisting of the sets \(\{x\}\) where \(x\) runs through E and the empty subset \(\varnothing\) of E; show that E' is a stable subset of \(\mathfrak{P}(E)\) under the law \((X,Y)\mapsto X\top Y\) (Exercise 1); deduce that, if \(\overline{E}\) denotes the set obtained by adjoining (Set Theory, II, § 4, no. 8) to E an element \(\omega\) , the law \(\top\) can be extended to \(\overline{E} \times \overline{E}\) , so that the law \(\top\) is identical with the law induced on E by this extended law.
\item
  Let \(\top\) be a law of composition not everywhere defined on E.
\end{enumerate}

\begin{enumerate}
\def\labelenumi{(\alph{enumi})}
\item
  For the law \((\mathbf{X}, \mathbf{Y}) \mapsto \mathbf{X} \top \mathbf{Y}\) (Exercise 1) between subsets of \(\mathbf{E}\) to be associative, it is necessary and sufficient that, for all \(x \in \mathbb{E}\) , \(y \in \mathbb{E}\) , \(z \in \mathbb{E}\) , if one of the two sides of the formula \((x \top y) \top z = x \top (y \top z)\) is defined, so is the other and is equal to it (use Exercise 1 to show that the condition is sufficient).
\item
  If this condition holds, show that Theorem 1 of no. 3 can be generalized as follows: if one of the two sides of formula (3) is defined, the other is defined and is equal to it.
\end{enumerate}

\begin{enumerate}
\def\labelenumi{\arabic{enumi}.}
\setcounter{enumi}{3}
\tightlist
\item
  \begin{enumerate}
  \def\labelenumii{(\alph{enumii})}
  \tightlist
  \item
    Given a set E, let \(\Phi\) be the set of mappings into E of any subset of E; if f, g, h are three elements of \(\Phi\) , show that if the composition \((f \circ g) \circ h\) is defined so is \(f \circ (g \circ h)\) , but that the converse is not true; if these two compositions are defined they are equal.
  \end{enumerate}
\end{enumerate}

\begin{enumerate}
\def\labelenumi{(\alph{enumi})}
\setcounter{enumi}{1}
\tightlist
\item
  Let \(\mathfrak{F}\) be a family of non-empty subsets of \(E\) such that no two have an element in common and let \(\Psi\) be the subset of \(\Phi\) consisting of the bijective mappings of a set of \(\mathfrak{F}\) onto a set of \(\mathfrak{F}\) . Show that, under the law induced by the law \(f \circ g\) on \(\Psi\) , the condition of Exercise 3(a) holds.
\end{enumerate}

\begin{enumerate}
\def\labelenumi{\arabic{enumi}.}
\setcounter{enumi}{4}
\item
  Show that the only triples \((m, n, p)\) of natural numbers \(\neq 0\) such that \((m^n)^p = m^{n^p}\) are: \((1, n, p)\) , \(n\) and \(p\) being arbitrary; \((m, n, 1)\) , \(m\) and \(n\) arbitrary and \((m, 2, 2)\) where \(m\) is arbitrary.
\item
  Let \(\top\) be a law of composition on a set E; let A be the subset of E consisting of the elements \(x\) such that \(x \top (y \top z) = (x \top y) \top z\) for all \(y \in E\) , \(z \in E\) ; show that A is a stable subset of E and that the law induced on A by \(\top\) is associative.
\item
  If \(\top\) is an associative law on \(\mathbf{E}\) and \(a\) and \(b\) two elements of \(\mathbf{E}\) , the sets \(\{a\} \top \mathbf{E}, \mathbf{E} \top \{b\}, \{a\} \top \mathbf{E} \top \{b\}, \mathbf{E} \top \{a\} \top \mathbf{E}\) are stable subsets of \(\mathbf{E}\) under the law \(\top\) .
\item
  Let \(\top\) be an associative law on E and \(a\) an element of E; for all \(x \in \mathrm{E}\) , \(y \in \mathrm{E}\) , we write \(x \perp y = x \top a \top y\) : show that the law \(\bot\) is associative.
\item
  On a set E the mappings \((x,y)\mapsto x\) and \((x,y)\mapsto y\) are opposite associative laws of composition.
\item
  Let X and Y be any two subsets of a set E and let \(X \top Y = X \cup Y\) if \(X \cap Y = \varnothing\) , \(X \top Y = E\) if \(X \cap Y \neq \varnothing\) ; show that the law of composition thus defined on \(\mathfrak{P}(E)\) is associative and commutative.
\item
  Let \(\top\) be an associative law on \(\mathbf{E}\) and \(\mathbf{A}\) and \(\mathbf{B}\) two subsets of \(\mathbf{E}\) which are stable under this law; show that, if \(\mathbf{B} \top \mathbf{A} \subset \mathbf{A} \top \mathbf{B}\) , \(\mathbf{A} \top \mathbf{B}\) is a stable subset of \(\mathbf{E}\) .
\item
  The only distinct natural numbers \(\neq 0\) which are permutable under the law \((x,y)\mapsto x^{y}\) are 2 and 4.
\item
  Show that under the law \((\mathbf{X}, \mathbf{Y}) \mapsto \mathbf{X} \circ \mathbf{Y}\) between subsets of \(E \times E\) the centre is the set consisting of \(\varnothing\) and the diagonal.
\item
  Show that under the law of composition \(f \circ g\) between mappings of E into E the centre consists of the identity mapping.
\item
  A law \(\top\) on a set \(E\) is called idempotent if all the elements of \(E\) are idempotent (no. 4) under this law, that is if \(x \top x = x\) for all \(x \in E\) . Show that, if a law \(\top\) on \(E\) is associative, commutative and idempotent, the relation \(x \top y = y\) is an order relation on \(E\) ; if we write \(x \leqslant y\) , any two elements \(x, y\) of \(E\) admit a least upper bound (under this order relation) equal to \(x \top y\) . Obtain the converse.
\item
  Let \(\top\) be an associative and commutative law of composition on a set E. Let \(n\) be an integer \(>0\) . The mapping \(x \mapsto \frac{n}{\top} x\) of E into E is a morphism.
\end{enumerate}

\BourbakiExercisesSection

\begin{enumerate}
\def\labelenumi{\arabic{enumi}.}
\item
  Let \(\top\) be a law of composition on a set E. Denoting by F the sum set (Set Theory, II, § 4, no. 8) of E and a set \(\{e\}\) consisting of a single element and identifying E and \(\{e\}\) with the corresponding subsets of F, show that it is possible in one and only one way to define on F a law of composition \(\overline{\top}\) which induces on E the law \(\top\) and under which \(e\) is identity element; if \(\top\) is associative, the law \(\overline{\top}\) is associative. (F is said to be derived from E ``by adjoining an identity element''.)
\item
  Let \(\mathsf{T}\) be a law everywhere defined on E.
\end{enumerate}

\begin{enumerate}
\def\labelenumi{(\alph{enumi})}
\item
  For \(\top\) to be associative, it is necessary and sufficient that every left translation \(\gamma_{x}\) be permutable with every right translation \(\delta_{y}\) in the set of mappings of E into E (with the law \(f \circ g\) ).
\item
  Suppose that E has an identity element. For T to be associative and commutative, it is necessary and sufficient that the mapping \((x,y)\mapsto x\top y\) be a morphism of the magma \(E\times E\) into the magma E.
\end{enumerate}

\begin{enumerate}
\def\labelenumi{\arabic{enumi}.}
\setcounter{enumi}{2}
\item
  In the set F of mappings of E into E, for the relation \(f \circ g = f \circ h\) to imply g = h, it is necessary and sufficient that f be injective; for the relation \(g \circ f = h \circ f\) to imply g = h, it is necessary and sufficient that f be surjective; for f to be cancellable (under the law \(\circ\) ), it is necessary and sufficient that f be bijective.
\item
  For there to exist on a set E a law of composition such that every permutation of E is an isomorphism of E onto itself under this law, it is necessary and sufficient that E have 0, 1 or 3 elements.
\item
  For \(2 \leqslant n \leqslant 5\) , determine on a set E with n elements all the laws everywhere defined admitting an identity element and under which all the elements of E are cancellable; for n = 5 show that there exist non-associative laws satisfying these conditions.
\end{enumerate}

(N.B. Exercises 6 to 17 inclusive refer to associative laws written multiplicatively; \(e\) denotes the identity when it exists, \(\gamma_a\) and \(\delta_a\) the translations by \(a \in E\) ; for every subset \(X\) of \(E\) we write \(\gamma_a(X) = aX\) , \(\delta_a(X) = Xa\) .)

\begin{enumerate}
\def\labelenumi{\arabic{enumi}.}
\setcounter{enumi}{5}
\item
  For an associative law on a finite set, every cancellable element is invertible (use no. 3, Proposition 3).
\item
  Given an associative law on a set E and an element \(x \in E\) , let A be the set of \(x^{n}\) for \(n \in N^{*}\) ; if there is an identity element, let B be the set of \(x^{n}\) for \(n \in N\) ; if further x is invertible, let C be the set of \(x^{a}\) for \(a \in Z\) . Show that, if A (resp.
\end{enumerate}

B, C) is infinite, it is isomorphic (with the law induced by the given law on E) to \(N^{*}\) (resp. N, Z) with addition.

\begin{enumerate}
\def\labelenumi{\arabic{enumi}.}
\setcounter{enumi}{7}
\item
  In the notation of Exercise 7, suppose A is finite; show that A contains one and only one idempotent (§ 1, no. 4) (observe that if \(x^p\) and \(x^q\) are idempotents, then \(x^p = x^{pq}\) , \(x^q = x^{pq}\) , hence \(x^p = x^q\) ); if \(h = x^p\) , the set of \(x^n\) for \(n \geqslant p\) is a stable subset D of E such that under the law induced on D, D is a group.
\item
  Under a multiplicative law on a set E let \(a\) be a left cancellable element of E.
\end{enumerate}

\begin{enumerate}
\def\labelenumi{(\alph{enumi})}
\item
  If there exists an element \(u\) such that \(au = a\) , show that \(ux = x\) for all \(x \in \mathbb{E}\) ; in particular, if \(xu = x\) for all \(x \in \mathbb{E}\) , \(u\) is identity element.
\item
  If there exist \(u \in E\) such that au = a and \(b \in E\) such that ab = u, show that ba = u (form aba); in particular, if there exists an identity element e and an element b such that ab = e, b is the inverse of a.
\end{enumerate}

\begin{enumerate}
\def\labelenumi{\arabic{enumi}.}
\setcounter{enumi}{9}
\tightlist
\item
  If a and b are two elements of E such that ba is left cancellable, show that a is left cancellable. Deduce that under an associative and commutative law on E the set S of non-cancellable elements of E is such that ES ⊆ S (and in particular is stable).
\end{enumerate}

¶ 11. E is called a left semi-group if every element x of E is left cancellable. (a) If u is an idempotent (§ 1, no. 4), then ux = x for all \(x \in E\) (use Exercise 9(a)); u is identity element for the law induced on Eu.

\begin{enumerate}
\def\labelenumi{(\alph{enumi})}
\setcounter{enumi}{1}
\item
  If \(u\) and \(v\) are two distinct idempotents of \(\mathbf{E}\) , then \(\mathrm{Eu} \cap \mathrm{Ev} = \varnothing\) and the sets \(\mathrm{Eu}\) and \(\mathrm{Ev}\) (with the laws induced by that on \(\mathbf{E}\) ) are isomorphic.
\item
  Let R be the complement of the union of the sets Eu, where u runs through the set of idempotents of E. Show that ER ⊂ R; it follows that R is a stable subset of E. If R is non-empty, then \(aR \neq R\) for all \(a \in R\) (use Exercise 9(a) to prove that, in the contrary case, R would contain an idempotent); in particular, R is then infinite. R is called the residue of the left semi-group E.
\item
  If R is non-empty and there exists at least one idempotent u in E (that is \(E \neq R\) ), for all \(x \in REu\) , \(xEu \neq Eu\) ; in particular, no element of REu is invertible in Eu and REu is an infinite set (use Exercise 8).
\item
  If E has an identity element e, e is the only idempotent in E and R is empty (note that E = Ee).
\item
  If there exists a right cancellable element \(a\) in E (in particular if the given law on E is commutative), either E has an identity element or \(\mathbf{E} = \mathbf{R}\) (note that if there exists an idempotent \(u\) , then \(xa = xua\) for all \(x \in \mathbf{E}\) ). For example, if E is the set of integers \(\geqslant 1\) , with addition, E is a commutative semi-group such that \(\mathbf{E} = \mathbf{R}\) .
\item
  If there exists \(a \in \mathbf{E}\) such that \(\delta_a\) is surjective, either \(\mathbf{E}\) has an identity element or \(\mathbf{E} = \mathbf{R}\) (examine separately the case where \(a \in \mathbf{R}\) and the case where \(a \in \mathbf{Eu}\) for an idempotent \(u\) ).
\end{enumerate}

¶ 12. Under a multiplicative law on E, let \(a \in E\) be such that \(\gamma_{a}\) is surjective.

\begin{enumerate}
\def\labelenumi{(\alph{enumi})}
\item
  Show that, if there exists \(u\) such that \(ua = a\) , then \(ux = x\) for all \(x \in E\) .
\item
  For an element \(b \in E\) to be such that ba is left cancellable, it is necessary and sufficient that \(\gamma_{a}\) be surjective and that b be left cancellable.
\end{enumerate}

¶ 13. Suppose that, for all \(x \in E\) , \(\gamma_x\) is surjective. Show that if, for an element \(a \in E\) , \(\gamma_a\) is a bijective mapping, \(\gamma_x\) is a bijective mapping for all \(x \in E\) (use Exercise 12(b)); this holds in particular if there exist two elements \(a, b\) of \(E\) such that \(ab = b\) (use Exercise 12(a)); \(E\) is then a left semi-group whose residue \(R\) is empty; moreover, for every idempotent \(u\) , every element of \(Eu\) is invertible in \(Eu\) .

¶ 14. For an associative law on a finite set E there exist minimal subsets of the form aE (that is minimal elements of the set of subsets of E of this form, ordered by inclusion).

\begin{enumerate}
\def\labelenumi{(\alph{enumi})}
\item
  If M = aE is minimal, then xM = xE = M for all \(x \in M\) ; with the law induced by that on E, M is a left semi-group (Exercise 11) in which every left translation is bijective (cf.~Exercise 13).
\item
  If \(\mathbf{M} = a\mathbf{E}\) and \(\mathbf{M}' = a'\mathbf{E}\) are minimal and distinct, then \(\mathbf{M} \cap \mathbf{M}' = \varnothing\) ; for all \(b \in \mathbf{M}\) , the mapping \(x' \mapsto bx'\) of \(\mathbf{M}'\) into \(\mathbf{M}\) is a bijective mapping of \(\mathbf{M}'\) onto \(\mathbf{M}\) . Deduce that there exist an idempotent \(u' \in \mathbf{M}'\) such that \(bu' = b\) and an idempotent \(u \in \mathbf{M}\) such that \(uu' = u\) (take \(u\) to be the idempotent such that \(bu = b\) ). Show that \(u'u = u'\) (consider \(uu'u\) ) and that every \(y' \in \mathbf{M}'\) such that \(y'u = y'\) belongs to \(\mathbf{M}'u'\) .
\item
  Show that the mapping \(x' \mapsto ux'\) of \(M'u'\) into \(M\) is an isomorphism of \(M'u'\) onto \(Mu\) ; deduce that \(M\) and \(M'\) are isomorphic left semi-groups.
\item
  Let \(M_{i}\) ( \(1 \leqslant i \leqslant r\) ) be the distinct minimal subsets of the form aE: deduce from (b) that the idempotents of each left semi-group \(M_{i}\) can be arranged in a sequence \((u_{ij})\) ( \(1 \leqslant j \leqslant s\) ) so that \(u_{ij}u_{kj} = u_{ij}\) for all i, j, k. If K is the union of the \(M_{i}\) , show that \(Eu_{ij} \subset K\) (note that, for all \(x \in E\) , \(xu_{ij}E\) is a minimal subset of the form aE); deduce that \(Eu_{ij}\) is the union of the \(u_{kj}Eu_{kj}\) for \(1 \leqslant k \leqslant r\) (show that \((\mathrm{Eu}_{ij}) \cap \mathbf{M}_{k} = u_{kj}\mathrm{Eu}_{kj}\) ) and that \(Eu_{ij}E = K\) . Prove finally that every minimal subset of the form Eb is identical with one of the s sets \(Eu_{ij}\) (note, using (a)), that \(Ebu_{ij}b = Eb\) and deduce that \(Eb \subset K\) ).
\end{enumerate}

\begin{enumerate}
\def\labelenumi{\arabic{enumi}.}
\setcounter{enumi}{14}
\tightlist
\item
  Let \((x_{i})_{1\leqslant i\leqslant n}\) be a finite sequence of left cancellable elements.
\end{enumerate}

\begin{enumerate}
\def\labelenumi{(\alph{enumi})}
\item
  Show that the relation \(x_{1}x_{2}\ldots x_{n} = e\) implies all the relations \(x_{i + 1}\ldots x_nx_1x_2\ldots x_i = e\) which are obtained from one another by ``cyclic permutation'' for \(1\leqslant i\leqslant n\) .
\item
  Deduce that if the composition of the sequence \((x_{i})\) is invertible each of the \(x_{i}\) is invertible.
\end{enumerate}

¶ 16. Let \(\top\) be an associative law on a set E and E* the set of cancellable elements of E; suppose that E* is non-empty and that every cancellable element is a central element. Let \(\mathfrak{F}\) denote the set of subsets \(\mathbf{X}\) of \(\mathbf{E}\) with the following property: there exists \(y \in \mathrm{E}^*\) such that \(\delta_y(\mathrm{E}) \subset \mathbf{X}\) .

\begin{enumerate}
\def\labelenumi{(\alph{enumi})}
\item
  Show that the intersection of two sets of \(\mathfrak{F}\) belong to \(\mathfrak{F}\) .
\item
  Let \(\Phi\) be the set of functions \(f\) defined on a set of \(\mathfrak{F}\) taking their values in
\end{enumerate}

E and such that, for \(X \in \mathfrak{F}\) , \(f(\mathbf{X})\) belongs to F. Let R denote the relation between elements f and g of \(\Phi\) : ``there exists a set \(X \in F\) such that the restrictions of f and g to X are identical''. Show that R is an equivalence relation; let \(\Psi = \Phi/R\) be the quotient set of \(\Phi\) under this relation.

\begin{enumerate}
\def\labelenumi{(\alph{enumi})}
\setcounter{enumi}{2}
\item
  Let \(f\) and \(g\) be two elements of \(\Phi\) , \(A \in \mathfrak{F}\) and \(B \in \mathfrak{F}\) the sets where \(f\) and \(g\) are respectively defined; there exists \(X \subset B\) and belonging to \(\mathfrak{F}\) such that \(g(X) \subset A\) ; if \(g_X\) is the restriction of \(g\) to \(X\) , show that the mapping \(f \circ g_X\) belongs to \(\Phi\) and that its class (mod. R) depends only on the classes of \(f\) and \(g\) (and not on \(X\) ); this class is called the composition of that of \(f\) and that of \(g\) ; show that the law of composition thus defined on \(\Psi\) is associative and has an identity element.
\item
  For all \(a \in \mathbb{E}\) , show that the left translation \(\gamma_{a}\) belongs to \(\Phi\) ; let \(\phi_{a}\) be its class (mod. R). Show that the mapping \(x \mapsto \phi_{x}\) is an isomorphism of E onto a submonoid of \(\Psi\) and that, if \(x \in \mathbb{E}^{*}\) , \(\phi_{x}\) is invertible in \(\Psi\) (consider the inverse mapping of \(\gamma_{x}\) , show that it belongs to \(\Phi\) and that its class (mod. R) is the inverse of \(\phi_{x}\) ). Deduce a generalization of Theorem 1 of no. 4.
\end{enumerate}

\begin{enumerate}
\def\labelenumi{\arabic{enumi}.}
\setcounter{enumi}{16}
\tightlist
\item
  Let E be a monoid and S a stable subset of E with the following properties:
\end{enumerate}

\((\alpha)\) For all \(s \in S\) and all \(a \in E\) , there exist \(t \in S\) and \(b \in E\) such that \(sb = at\) .

(β) For all \(a, b \in \mathbb{E}\) and all \(s \in \mathbb{S}\) such that \(sa = sb\) , there exist \(t \in \mathbb{S}\) such that \(at = bt\) .

\begin{enumerate}
\def\labelenumi{(\alph{enumi})}
\tightlist
\item
  In E × S let \((a, s) \sim (b, t)\) denote the relation: ``there exist \(s'\) and \(t'\) in S such that \(tt' = ss'\) and \(as' = bt'\) .''
\end{enumerate}

Show that \(\sim\) is an equivalence relation. We write \(\overline{\mathbf{E}} = \mathbf{E} \times \mathbf{S}/\sim\) , denote by \(as^{-1}\) the equivalence class of \((a, s) (\text{mod } \sim)\) and by \(\varepsilon: \mathbf{E} \to \overline{\mathbf{E}}\) the mapping \(a \mapsto ae^{-1}\) .

\begin{enumerate}
\def\labelenumi{(\alph{enumi})}
\setcounter{enumi}{1}
\item
  Show that there exists on \(\overline{\mathbf{E}}\) one and only one monoid structure such that \(\varepsilon\) is a unital homomorphism and such that, for all \(s \in \mathbf{S}\) , \(\varepsilon(s)\) is invertible.
\item
  Show that \((\overline{\mathbf{E}},\varepsilon)\) has the universal property described in Theorem 1 (no. 4).
\item
  Show that \(\varepsilon\) is injective if and only if the elements of S are cancellable.
\end{enumerate}

\BourbakiExercisesSection

\begin{enumerate}
\def\labelenumi{\arabic{enumi}.}
\tightlist
\item
  Let \(\alpha \mapsto f_{\alpha}\) be an action of \(\Omega\) on E. Let F be the image of \(\Omega\) in \(\mathbf{E}^{\mathbb{E}}\) and G the stable subset of \(\mathbf{E}^{\mathbb{E}}\) under the law \((f,g) \mapsto f \circ g\) generated by F and the identity mapping of E.
\end{enumerate}

\begin{enumerate}
\def\labelenumi{(\alph{enumi})}
\item
  Show that every subset of E which is stable under the action of \(\Omega\) is also stable under the restriction to G of the canonical action of \(E^{E}\) on E (no. 1, Example 3).
\item
  Let X be a subset of E; show that the stable subset of E generated by X is the set of \(f(x)\) where f runs through G and x runs through X.
\end{enumerate}

\begin{enumerate}
\def\labelenumi{\arabic{enumi}.}
\setcounter{enumi}{1}
\item
  Let \(\bot\) be an internal law on E which is doubly distributive with respect to the associative internal law \(\top\) ; show that, if \(x \perp x'\) and \(y \perp y'\) are cancellable under the law \(\top\) , \(x \perp y'\) is permutable with \(y \perp x'\) under the law \(\top\) (calculate the composition \((x \top y) \perp (x' \top y')\) in two different ways). In particular, if the law \(\bot\) has an identity element, two cancellable elements under the law \(\top\) are permutable under this law; if all the elements of E are cancellable under the law \(\top\) , this law is commutative.
\item
  Let \(\top\) and \(\bot\) be two internal laws on a set \(E\) such that \(\bot\) is right distributive with respect to \(\top\) .
\end{enumerate}

\begin{enumerate}
\def\labelenumi{(\alph{enumi})}
\item
  If the law \(\top\) has an identity element \(e\) , \(x \perp e\) is idempotent under the law \(\top\) for all \(x \in \mathbb{E}\) ; if further there exists \(y\) such that \(x \perp y\) is cancellable under the law \(\top\) , then \(x \perp e = e\) .
\item
  If the law \(\perp\) has an identity element u and there exists \(z \in E\) which is cancellable for both laws \(\perp\) and \(\top\) , u is cancellable under the law \(\top\) .
\end{enumerate}

¶ 4. Let \(\top\) and \(\bot\) be two internal laws on E, each with an identity element; if the left action of E on itself derived from each of the laws is distributive with respect to the other, every element of E is idempotent under both these laws (if \(e\) is the identity element for \(\top\) , \(u\) the identity element for \(\bot\) , prove first that \(e \perp e = e\) , noting that \(e = e \perp (u \top e)\) ).

\begin{enumerate}
\def\labelenumi{\arabic{enumi}.}
\setcounter{enumi}{4}
\tightlist
\item
  On a set E suppose three internal laws of composition are given: an addition (not necessarily associative nor commutative), a multiplication (not necessarily associative) and a law denoted by \(\top\) . Suppose that the multiplication has an identity element \(e\) , that the law \(\top\) is right distributive with respect to multiplication and the law \(\top\) is left distributive with respect to addition. Show that, if there exist \(x, y, z\) such that \(x \top z, y \top z\) and \((x + y) \top z\) are cancellable under multiplication, then \(e + e = e\) (use Exercise 3(a)).
\end{enumerate}

¶ 6. An internal law of composition \(\top\) on E is said to determine on E a quasi-group structure if, for all \(x \in E\) , the left and right translations \(\gamma_x\) and \(\delta_x\) are bijective mappings of E onto itself. A quasi-group is said to be distributive if the law \(\top\) is doubly distributive with respect to itself.

\begin{enumerate}
\def\labelenumi{(\alph{enumi})}
\item
  Determine all the distributive quasi-group structures on a set of \(n\) elements for \(2 \leqslant n \leqslant 6\) .
\item
  *Show that the set \(\mathbf{Q}\) of rational numbers, with the law of composition \((x,y)\mapsto\frac{1}{2}(x+y)\) , is a distributive quasi-group.*
\item
  In a distributive quasi-group E, every element is idempotent. Deduce that, if E has more than one element, the law \(\top\) can neither have an identity element nor be associative.
\item
  The right and left translations of a distributive quasi-group E are automorphisms of E.
\item
  If E is a finite distributive quasi-group, the structure induced on every stable subset of E is a distributive quasi-group structure.
\item
  Let E be a distributive quasi-group. If R is an equivalence relation which is left (resp. right) compatible (no. 3) with the law T, the classes mod. R are stable subsets of E. If E is finite, all these classes are obtained one from another by left (resp. right) translation; under the same conditions, if R is compatible with the law T, the quotient structure on E/R is a distributive quasi-group structure.
\item
  Let E be a distributive quasi-group. The set \(A_{a}\) of elements of E which are permutable with a given element a is stable; if E is finite, for all \(x \in A_{a}\) , \(A_{x} = A_{a}\) (note, using (e), that there exists \(y \in A_{a}\) such that \(x = a \top y\) ) and when x runs through E the sets \(A_{x}\) are identical with the equivalence classes with respect to a relation compatible with the law T.
\item
  If E is a finite distributive quasi-group and the law \(\top\) is commutative, the number of elements of E is odd (consider the ordered pairs \((x,y)\) of elements of E such that \(x\top y=y\top x=a\) , where a is given).
\end{enumerate}

\begin{enumerate}
\def\labelenumi{\arabic{enumi}.}
\setcounter{enumi}{6}
\tightlist
\item
  Suppose given on a set E an (associative and commutative) addition under which all the elements of E are invertible and an (associative) multiplication which is doubly distributive with respect to addition \(*\) (in other words, a ring structure) \(_{*}\) ; write \([x,y]=xy-yx\) ; the law \((x,y)\mapsto[x,y]\) is doubly distributive with respect to addition. For x and y to be permutable under multiplication, it is necessary and sufficient that \([x,y]=0\) ; and we have the identities
\end{enumerate}

\[
[ x, y ] = - [ y, x ]; \quad [ x, [ y, z ] ] + [ y, [ z, x ] ] + [ z, [ x, y ] ] = 0
\]

(the second is known as the ``Jacobi identity''). The second identity may also be written

\[
[ x, [ y, z ] ] - [ [ x, y ], z ] = [ [ z, x ], y ]
\]

which expresses the ``deviation from associativity'' of the law \([x, y]\) .

\begin{enumerate}
\def\labelenumi{\arabic{enumi}.}
\setcounter{enumi}{7}
\tightlist
\item
  With the same hypotheses as in Exercise 7, let \(x \top y = xy + yx\) ; then the law \(\top\) is commutative and doubly distributive with respect to addition but not in general associative.
\end{enumerate}

\begin{enumerate}
\def\labelenumi{(\alph{enumi})}
\item
  For all \(x \in \mathbb{E}\) , show that \(\top^{m+n} x = \left( \begin{array}{c} m \\ \top \end{array} x \right) \top \left( \begin{array}{c} n \\ \top \end{array} x \right)\) .
\item
  If we write \([x, y, z] = (x \top y) \top z - x \top (y \top z)\) (deviation from associativity of the law \(\top\) ), prove the identities
\end{enumerate}

\[
[ x, y, z ] + [ z, y, x ] = 0
\]

\[
[ x, y, z ] + [ y, z, x ] + [ z, x, y ] = 0
\]

\[
[ x \top y, u, z ] = [ u, [ (x \top y), z) ] ]
\]

(the notation \([x, y]\) meaning the same as in Exercise 7)

\[
[ x \top y, u, z ] + [ y \top z, u, x ] + [ z \top x, u, y ] = 0.
\]

¶ 9. Suppose given on E an (associative and commutative) addition under which all the elements of E are invertible and a multiplication which is non-associative, but commutative and doubly distributive with respect to addition. Suppose further that \(n \in \mathbf{Z}\) , \(n \neq 0\) and \(nx = 0\) imply \(x = 0\) in E. Show that if, writing \([x, y, z] = (xy)z - x(yz)\) , the identity

\[
[ x y, u, z ] + [ y z, u, z ] + [ z x, u, y ] = 0
\]

holds, then \(x^{m + n} = x^m x^n\) for all \(x\) (show, by induction on \(p\) , that the identity \([x^q, y, x^{p - q}] = 0\) holds for \(1 \leqslant q < p\) ).

\begin{enumerate}
\def\labelenumi{\arabic{enumi}.}
\setcounter{enumi}{9}
\tightlist
\item
  Let E be a commutative monoid whose law is written additively and whose identity element is denoted by 0. Let \(\top\) be an internal law on E which is distributive with respect to addition and such that \(0 \top x = x \top 0 = 0\) for all \(x \in \mathbb{E}\) . Let S be a subset of E which is stable under addition and under each of the external laws derived from \(\top\) ; let \(\overline{\mathbb{E}}\) denote the monoid of differences of E with respect to S under addition and \(\varepsilon\) the canonical homomorphism of E into \(\overline{\mathbb{E}}\) . Then there exists on \(\overline{\mathbb{E}}\) one and only one law \(\overline{\top}\) which is distributive with respect to addition and such that \(\varepsilon\) is a homomorphism for the laws \(\top\) and \(\overline{\top}\) . If the law \(\top\) is associative (resp. commutative), so is the law \(\overline{\top}\) .
\end{enumerate}

\BourbakiExercisesSection

\begin{enumerate}
\def\labelenumi{\arabic{enumi}.}
\tightlist
\item
  Determine all the group structures on a set of \(n\) elements for \(2 \leqslant n \leqslant 6\) (cf.~§2, Exercise 5). Determine the subgroups and quotient groups of these groups and also their Jordan-Hölder series.
\end{enumerate}

¶ 2. (a) An associative law \((x,y)\mapsto xy\) on a set E is a group law if there exists \(e\in\mathrm{E}\) such that, for all \(x\in\mathrm{E}\) , \(ex=x\) and if, for all \(x\in\mathrm{E}\) , there exists \(x'\) such that \(x^{\prime}x=e\) (show that \(xx^{\prime}=e\) by considering the composition \(x^{\prime}xx^{\prime}\) ; deduce that \(e\) is the identity element).

\begin{enumerate}
\def\labelenumi{(\alph{enumi})}
\setcounter{enumi}{1}
\tightlist
\item
  Show that the same is true if, for all \(x \in E\) , the left translation \(\gamma_{x}\) is a mapping of E onto E and if there exists some \(a \in E\) such that the right translation \(\delta_{a}\) is a mapping of E onto E.
\end{enumerate}

\begin{enumerate}
\def\labelenumi{\arabic{enumi}.}
\setcounter{enumi}{2}
\item
  In a group G every finite non-empty stable subset H is a subgroup of G (cf.~§ 2, Exercise 6).
\item
  Let A and B be two subgroups of a group G.
\end{enumerate}

\begin{enumerate}
\def\labelenumi{(\alph{enumi})}
\item
  Show that the least subgroup containing A and B (that is the subgroup generated by \(\mathbf{A} \cup \mathbf{B}\) ) is identical with the set of compositions of the sequences \((x_{i})_{1 \leqslant i \leqslant 2n + 1}\) of an (arbitrary) odd number of elements such that \(x_{i} \in \mathbf{A}\) for \(i\) odd and \(x_{i} \in \mathbf{B}\) for \(i\) even.
\item
  For AB to be a subgroup of G (in which case it is the subgroup generated by A ∪ B), it is necessary and sufficient that A and B be permutable, that is that AB = BA.
\item
  If A and B are permutable and C is a subgroup containing A, A is permutable with \(B \cap C\) and \(\mathbf{A}(\mathbf{B} \cap \mathbf{C}) = \mathbf{C} \cap (\mathbf{AB})\) .
\end{enumerate}

\begin{enumerate}
\def\labelenumi{\arabic{enumi}.}
\setcounter{enumi}{4}
\item
  If a subgroup of a group G has index 2 it is normal in G.
\item
  Let \((\mathbf{G}_{\alpha})\) be a family of normal subgroups of a group \(\mathbf{G}\) such that \(\bigcap_{\alpha} \mathbf{G}_{\alpha} = \{e\}\) ; show that \(\mathbf{G}\) is isomorphic to a subgroup of the product group \(\prod_{\alpha} (\mathbf{G} / \mathbf{G}_{\alpha})\) .
\item
  If G is the direct product of two subgroups A and B and H is a subgroup of G such that \(A \subset H\) , H is the direct product of A and \(H \cap B\) .
\item
  Let A and A' be two groups and G a subgroup of A × A'. We write:
\end{enumerate}

\[
\mathrm{N} = \mathrm{G} \cap (\mathrm{A} \times \{e \}), \quad \mathrm{H} = \operatorname{pr} _ {1} (\mathrm{G})
\]

\[
\mathrm{N} ^ {\prime} = \mathrm{G} \cap (\{e \} \times \mathrm{A} ^ {\prime}), \quad \mathrm{H} ^ {\prime} = \operatorname{pr} _ {2} (\mathrm{G}).
\]

\begin{enumerate}
\def\labelenumi{(\alph{enumi})}
\tightlist
\item
  Show that N is normal in H and N' is normal in H'; define isomorphisms
\end{enumerate}

\[
\mathrm{H} / \mathrm{N} \rightarrow \mathrm{G} / (\mathrm{N} \times \mathrm{N} ^ {\prime}) \rightarrow \mathrm{H} ^ {\prime} / \mathrm{N} ^ {\prime}.
\]

\begin{enumerate}
\def\labelenumi{(\alph{enumi})}
\setcounter{enumi}{1}
\tightlist
\item
  Suppose that H = A, \(H' = A'\) , that the groups A and \(A'\) are of finite length and that no quotient of a Jordan-Hölder series of A is isomorphic to a quotient of a Jordan-Hölder series of \(A'\) . Show that \(G = A \times A'\) .
\end{enumerate}

\begin{enumerate}
\def\labelenumi{\arabic{enumi}.}
\setcounter{enumi}{8}
\item
  Let G be a group which does not consist only of the identity element. Suppose that there exists a finite subset S of G such that G is generated by S (resp. by the elements \(gsg^{-1}, g \in G, s \in S\) ). Let N be the set of subgroups of G (resp. of normal subgroups of G) distinct from G. Show that N, ordered by inclusion, is inductive. Deduce the existence of a normal subgroup H of G such that G/H is simple.
\item
  Let H be a normal subgroup of a group G contained in the centre of G. Show that if G/H is a monogenous group G is commutative.
\item
  If all the elements of a group G other than the identity element are of order 2, G is commutative; if G is finite, its order n is a power of 2 (argue by induction on n).
\item
  Let G be a group such that, for a fixed integer \(n > 1\) , \((xy)^n = x^n y^n\) for all \(x \in G\) , \(y \in G\) . If \(G^{(n)}\) denotes the set of \(x^n\) , where \(x\) runs through G, and \(G_{(n)}\) the set of \(x \in G\) such that \(x^n = e\) , show that \(G^{(n)}\) and \(G_{(n)}\) are normal subgroups of G; if G is finite, the order of \(G^{(n)}\) is equal to the index of \(G_{(n)}\) . Show that, for all \(x, y\) in G, also \(x^{1-n}y^{1-n} = (xy)^{1-n}\) and deduce that \(x^{n-1}y^n = y^n x^{n-1}\) ; conclude from this that the set of elements of G of the form \(x^{n(n-1)}\) generates a commutative subgroup of G.
\item
  Let G be a group. If S is a simple group, S is said to occur in G if there exist two subgroups H, H' of G with H a normal subgroup of H' such that H'/H is isomorphic to S. Let in(G) denote the set of isomorphism classes of simple groups occurring in G.
\end{enumerate}

\begin{enumerate}
\def\labelenumi{(\alph{enumi})}
\item
  Show that in(G) = ∅ ⇔ G = \{e\}.
\item
  If \(\mathbf{H}\) is a subgroup of \(\mathbf{G}\) , show that \(\operatorname{in}(\mathbf{H}) \subset \operatorname{in}(\mathbf{G})\) ; if moreover \(\mathbf{H}\) is normal, then
\end{enumerate}

\[
\operatorname{in} (\mathrm{G}) = \operatorname{in} (\mathrm{H}) \cup \operatorname{in} (\mathrm{G} / \mathrm{H}).
\]

\begin{enumerate}
\def\labelenumi{(\alph{enumi})}
\setcounter{enumi}{2}
\item
  Let \(\mathbf{G}_1\) and \(\mathbf{G}_2\) be two groups. Show the equivalence of the two following properties:
\item
  \(\operatorname{in}(\mathbf{G}_1) \cap \operatorname{in}(\mathbf{G}_2) = \varnothing\) ;
\end{enumerate}

\begin{enumerate}
\def\labelenumi{(\roman{enumi})}
\setcounter{enumi}{1}
\tightlist
\item
  Every subgroup of \(\mathbf{G}_1\times \mathbf{G}_2\) is of the form \(\mathbf{H}_1\times \mathbf{H}_2\) with \(\mathbf{H}_1\subset \mathbf{G}_1\) and \(\mathbf{H}_2\subset \mathbf{G}_2\) .
\end{enumerate}

(Use Exercise 8.)

*When \(\mathbf{G}_1\) and \(\mathbf{G}_2\) are finite groups, show that (i) and (ii) are equivalent to:

\begin{enumerate}
\def\labelenumi{(\roman{enumi})}
\setcounter{enumi}{2}
\tightlist
\item
  \(\operatorname{Card}(\mathbf{G}_1)\) and \(\operatorname{Card}(\mathbf{G}_2)\) are relatively prime.*
\end{enumerate}

\begin{enumerate}
\def\labelenumi{\arabic{enumi}.}
\setcounter{enumi}{13}
\tightlist
\item
  Let G be a group and H a subgroup of G. Any subset T of G which meets every left coset mod. H in one and only one point is called a representative system of the left cosets mod. H (cf.~Set Theory, II, § 6, no. 2); this condition is equivalent to the following:
\end{enumerate}

The mapping \((x,y)\mapsto xy\) is a bijection of \(\mathbf{T}\times \mathbf{H}\) onto \(\mathbf{G}\) .

\begin{enumerate}
\def\labelenumi{(\alph{enumi})}
\item
  Let \(\mathbf{T}\) be such a representative system; for all \(g \in \mathbf{G}\) and all \(t \in \mathbf{T}\) , let \(x(g, t) \in \mathbf{T}\) and \(y(g, t) \in \mathbf{H}\) be such that \(gt = x(g, t)y(g, t)\) . Let \(\mathbf{S}\) be a subset of \(\mathbf{G}\) generating \(\mathbf{G}\) . Show that the elements \(y(g, t), g \in \mathbf{S}, t \in \mathbf{T}\) generate \(\mathbf{H}\) . (If \(\mathbf{H}'\) denotes the subgroup of \(\mathbf{H}\) generated by these elements, show that \(\mathbf{T}.\mathbf{H}'\) is stable under the \(\boldsymbol{\gamma}_g, g \in \mathbf{S}\) , hence also under the \(\boldsymbol{\gamma}_g, g \in \mathbf{G}\) , and deduce that \(\mathbf{T}.\mathbf{H}' = \mathbf{G}\) , whence \(\mathbf{H}' = \mathbf{H}\) .)
\item
  Suppose that (G:H) is finite. Show that G can be generated by a finite subset if and only if H can.
\end{enumerate}

¶ 15. Let F be a set of stable subgroups of a group with operators G; F is said to satisfy the maximal condition (resp. minimal condition) if every subset of F, ordered by inclusion, has a maximal (resp. minimal) element.

Suppose that the set of all stable subgroups of a group G satisfies the minimal condition.

\begin{enumerate}
\def\labelenumi{(\alph{enumi})}
\item
  Prove that there exists no stable subgroup of G isomorphic to G and distinct from G (argue by reductio ad absurdum by showing that the hypothesis would imply the existence of a strictly decreasing infinite sequence of stable subgroups of G).
\item
  The minimal elements of the set of normal stable subgroups of G not consisting only of e are called minimal normal stable subgroups. Let M be a set of minimal normal stable subgroups of G and S the smallest stable subgroup of G containing all the subgroups belonging to M; show that S is the direct product of a finite number of minimal normal stable subgroups of G (let \((\mathbf{M}_{n})\) be a sequence of minimal normal stable subgroups of G belonging to M and such that \(M_{n+1}\) is not contained in the stable subgroup generated by the union of \(M_{1}, M_{2}, \ldots, M_{n}\) ; let \(S_{k}\) be the stable subgroup generated by the union of the \(M_{n}\) of index \(n \geqslant k\) ; show that \(S_{k+1} = S_{k}\) after a certain rank and therefore that \((\mathbf{M}_{n})\) is a finite sequence; finally use no. 9, Proposition 15).
\item
  If G is a group without operators, show that every minimal normal subgroup M of G is the direct product of a finite number of simple groups isomorphic to one another (let N be a minimal normal subgroup of M; show that M is the smallest subgroup of G containing all the subgroups \(aNa^{-1}\) where a runs through G, and then apply (b) to the group M).
\end{enumerate}

\begin{enumerate}
\def\labelenumi{\arabic{enumi}.}
\setcounter{enumi}{15}
\tightlist
\item
  If the set of stable subgroups of a group with operators G satisfies the maximal and minimal conditions (Exercise 15), G has a Jordan-Hölder series (consider, for a subgroup H of G, a maximal element of the set of normal stable subgroups of H distinct from H).
\end{enumerate}

¶ 17. Let G be a group with operators; a composition series (G \(_{i}\) ) of G is called normal if all the G \(_{i}\) are normal stable subgroups of G; a normal series Σ is called principal if it is strictly decreasing and there exists no normal series distinct from Σ, finer than Σ and strictly decreasing.

\begin{enumerate}
\def\labelenumi{(\alph{enumi})}
\item
  If \((\mathbf{G}_{i})\) and \((\mathbf{H}_{j})\) are two normal series of G, show that there exist two equivalent normal series finer respectively than \((\mathbf{G}_{i})\) and \((\mathbf{H}_{j})\) (apply Schreier's Theorem by considering a suitable domain of operators for G). Give a second proof of this proposition by ``inserting'' the subgroups \(G_{ij}^{\prime} = G_{i} \cap (G_{i+1}H_{j})\) and \(H_{ji}^{\prime} = H_{j} \cap (H_{j+1}G_{i})\) in the sequences \((\mathbf{G}_{i})\) and \((\mathbf{H}_{j})\) respectively.
\item
  If G has a principal series, any two principal series of G are equivalent; for every strictly decreasing normal series \(\Sigma\) there exists a principal series finer than \(\Sigma\) . Deduce that, for G to possess a principal series, it is necessary and sufficient that the set of normal stable subgroups of G satisfy the maximal and minimal conditions.
\item
  If G is a group without operators and possesses a principal series and if the set of subgroups of G satisfies the minimal condition, every quotient group
\end{enumerate}

\(G_{i}/G_{i+1}\) is the direct product of a finite number of simple subgroups isomorphic to one another (use Exercise 15).

¶ 18. (a) Let G be a group with operators generated by the union of a family \((\mathrm{H}_{i})_{i\in\mathbf{I}}\) of simple normal stable subgroups of G. Show that there exists a subset J of I such that G is the restricted sum of the family \((\mathrm{H}_{i})_{i\in\mathbf{J}}\) (apply Zorn's Lemma to the set of subsets K of I such that the subgroup generated by the union of the \(H_{i}\) for \(i\in K\) is the restricted sum of this family).

\begin{enumerate}
\def\labelenumi{(\alph{enumi})}
\setcounter{enumi}{1}
\tightlist
\item
  Let A be a normal stable subgroup of G. Show that there exists a subset \(J_{A}\) of I such that G is the restricted sum of A and the \(H_{i}\) for \(i \in J_{A}\) . Deduce that A is isomorphic to the restricted sum of a subfamily of the \(H_{i}\) .
\end{enumerate}

¶ 19. (a) Let G be a group such that every normal subgroup of G distinct from G is contained in a direct factor of G distinct from G. Show that G is the restricted sum of a family of simple subgroups. (Let K be the subgroup of G generated by the union of all the simple subgroups of G. Suppose K ≠ G; if x \ensuremath{\in} G = K, consider a normal subgroup M of G which is maximal amongst those containing K and not containing x. If G = M' × S, where M' contains M and S ≠ \{e\}, show that x ∉ M' and therefore M' = M; on the other hand, if N is a normal subgroup of S, show that M × N = G and conclude that S is simple, whence a contradiction. Conclude with the aid of Exercise 18.) Obtain the converse (cf.~Exercise 18).

\begin{enumerate}
\def\labelenumi{(\alph{enumi})}
\setcounter{enumi}{1}
\tightlist
\item
  In order that there exist in a group G a family \((\mathrm{N}_{\alpha})\) of normal simple subgroups such that the \(G/N_{\alpha}\) are simple and \(\bigcap_{\alpha} N_{\alpha} = \{e\}\) , it is necessary and sufficient that, for every normal subgroup \(N \neq \{e\}\) of G, there exist a normal subgroup \(N' \neq G\) such that \(NN' = G\) . (To see that the condition is necessary, consider an \(x \in N\) distinct from e and an \(N_{\alpha}\) not containing x. To see that the condition is sufficient, for all \(x \neq e\) in G, consider the normal subgroup N generated by x and a normal subgroup \(N' \neq G\) such that \(NN' = G\) . Show that if a normal subgroup M of G is maximal amongst the normal subgroups of G containing \(N'\) and not containing x, it is maximal in the set of all normal subgroups \(\neq G\) and deduce the conclusion.)
\end{enumerate}

\begin{enumerate}
\def\labelenumi{\arabic{enumi}.}
\setcounter{enumi}{19}
\tightlist
\item
  \begin{enumerate}
  \def\labelenumii{(\alph{enumii})}
  \tightlist
  \item
    Let G be a group, the direct product of two subgroups A and B. Let \(C_{A}\) be the centre of A and let N be a normal subgroup of G such that \(N \cap A = \{1\}\) . Show that \(N \subset C_{A} \times B\) .
  \end{enumerate}
\end{enumerate}

\begin{enumerate}
\def\labelenumi{(\alph{enumi})}
\setcounter{enumi}{1}
\tightlist
\item
  Let \((\mathbf{S}_i)_{i\in I}\) be a finite family of non-commutative simple groups. Show that every normal subgroup of \(\prod_{i\in I}S_i\) is equal to one of the partial products \(\prod_{i\in J}S_i\) where \(J\) is a subset of \(I\) .
\end{enumerate}

\begin{enumerate}
\def\labelenumi{\arabic{enumi}.}
\setcounter{enumi}{20}
\tightlist
\item
  Let G be a group of finite length. Show the equivalence of the following properties:
\end{enumerate}

\begin{enumerate}
\def\labelenumi{(\alph{enumi})}
\item
  G is a product of simple groups isomorphic to one another.
\item
  No subgroup of G distinct from \(\{1\}\) and G is stable under all the automorphisms of G.
\end{enumerate}

¶ 22. Let E be a quasi-group (§ 3, Exercise 6) written multiplicatively, possessing an idempotent e and such that \((xy)(zt) = (xz)(yt)\) for all x, y, z, t. Let \(u(x)\) denote the element of E such that \(u(x)e = x\) and \(v(x)\) the element of E such that \(ev(x) = x\) ; show that the law of composition \((x, y) \mapsto u(x)v(y)\) is a commutative group law on E under which e is identity element, that the mappings \(x \mapsto xe\) and \(y \mapsto ey\) are permutable endomorphisms of this group structure and that \(xy = u(xe)v(ey)\) (start by establishing the identities \(e(xy) = (ex)(ey)\) , \((xy)e = (xe)(ye)\) , \(e(xe) = (ex)e\) ; then note that the relations x = y, ex = ey and xe = ye are equivalent). Obtain the converse.

¶ 23. Consider on a set E an internal law, not everywhere defined, written multiplicatively and satisfying the following conditions:

\begin{enumerate}
\def\labelenumi{(\arabic{enumi})}
\item
  if one of the compositions \((xy)z, x(yz)\) is defined, so is the other and they are equal;
\item
  if \(x, x', y\) are such that \(xy\) and \(x'y\) (resp. \(yx\) and \(yx'\) ) are defined and equal, then \(x = x'\) ;
\item
  for all \(x \in \mathbf{E}\) , there exist three elements \(e_x, e_x'\) and \(x^{-1}\) such that \(e_xx = x\) , \(xe_x' = x\) , \(x^{-1}x = e_x'\) ; \(e_x\) is called the left unit of \(x\) , \(e_x'\) the right unit of \(x\) , \(x^{-1}\) the inverse of \(x\) (by an abuse of language).
\end{enumerate}

\begin{enumerate}
\def\labelenumi{(\alph{enumi})}
\item
  Show that the compositions \(xx^{-1}, x^{-1}e_x, e_x'x^{-1}, e_x e_x, e_x' e_x'\) are defined and that \(xx^{-1} = e_x, x^{-1}e_x = e_x'x^{-1} = x^{-1}, e_x e_x = e_x, e_x' e_x' = e_x'\) .
\item
  Every idempotent \(e\) of \(\mathbf{E}\) (§ 1, no. 4) is a left unit for all the \(x\) such that \(ex\) is defined and right unit for all the \(y\) such that \(ye\) is defined.
\item
  For the composition \(xy\) to be defined, it is necessary and sufficient that the right unit of \(x\) be the same as the left unit of \(y\) (to see that the condition is sufficient, use the relation \(e_y = yy^{-1}\) ); if \(xy = z\) , then \(x^{-1}z = y\) , \(zy^{-1} = x\) , \(y^{-1}x^{-1} = z^{-1}\) , \(z^{-1}x = y^{-1}\) , \(yz^{-1} = x^{-1}\) (the compositions appearing on the left-hand sides of these relations being defined).
\item
  For any two idempotents \(e, e'\) of \(\mathbf{E}\) , let \(\mathbf{G}_{e, e'}\) denote the set of \(x \in \mathbf{E}\) such that \(e\) is left unit and \(e'\) is right unit of \(x\) . Show that, if \(\mathbf{G}_{e, e}\) is non-empty, it is a group under the law induced by that on \(\mathbf{E}\) .
\end{enumerate}

E is called a groupoid if it also satisfies the following condition:

\begin{enumerate}
\def\labelenumi{(\arabic{enumi})}
\setcounter{enumi}{3}
\tightlist
\item
  for all idempotents \(e, e', G_{e, e'}\) is non-empty.
\end{enumerate}

\begin{enumerate}
\def\labelenumi{(\alph{enumi})}
\setcounter{enumi}{4}
\item
  In a groupoid E, if \(a \in G_{e,e'}\) , show that \(x \mapsto xa\) is a bijective mapping of \(G_{e,e}\) onto \(G_{e,e'}, y \mapsto ay\) a bijective mapping of \(G_{e',e'}\) onto \(G_{e,e'}\) and \(x \mapsto a^{-1}xa\) an isomorphism of the group \(G_{e,e}\) onto the group \(G_{e',e'}\) .
\item
  Show that the law defined in § 1, Exercise 4(b) determines a groupoid structure on the set \(\Psi\) , if all the sets of the family \(\mathfrak{F}\) are equipotent.
\end{enumerate}

\begin{enumerate}
\def\labelenumi{\arabic{enumi}.}
\setcounter{enumi}{23}
\tightlist
\item
  \begin{enumerate}
  \def\labelenumii{(\alph{enumii})}
  \tightlist
  \item
    Given any set E, consider in the set \(E \times E\) the law of composition, not everywhere defined, written multiplicatively, under which the composition of \((x, y)\) and \((y', z)\) is only defined if \(y' = y\) and in this case has the value \((x, z)\) . Show that \(\mathbf{E} \times \mathbf{E}\) , with this law of composition, is a groupoid (Exercise 23).
  \end{enumerate}
\end{enumerate}

\begin{enumerate}
\def\labelenumi{(\alph{enumi})}
\setcounter{enumi}{1}
\tightlist
\item
  Let \((x, y) \equiv (x', y')\) (R) be an equivalence relation compatible with the above composition (that is such that \((x, y) \equiv (x', y')\) and \((y, z) \equiv (y', z')\) imply \((x, z) \equiv (x', z'))\) ; suppose further that the relation R satisfies the following condition: for all \(x, y, z\) there exists one and only one \(t \in \mathbf{E}\) such that \((x, y) \equiv (z, t)\) and there exists a \(u \in \mathbf{E}\) such that \((x, y) \equiv (u, z)\) .
\end{enumerate}

Show that under these conditions the quotient structure by R of the groupoid structure on \(E \times E\) is a group structure (prove first that the quotient law is everywhere defined, since, if \(\dot{x}, \dot{y}, \dot{z}\) are three classes such that \(\dot{x}\dot{y} = \dot{x}\dot{z}\) , then \(\dot{y} = \dot{z}\) ; establish finally that, for all \(x \in E, y \in E, (x, x) \equiv (y, y)\) ).

\begin{enumerate}
\def\labelenumi{(\alph{enumi})}
\setcounter{enumi}{2}
\tightlist
\item
  Let G be the group obtained by giving the quotient of \(E \times E\) by R the above structure. Let a be any element of E; if, for all \(x \in E\) , \(f_{a}(x)\) denotes the class mod. R of \((a, x)\) , show that \(f_{a}\) is a bijective mapping of E onto G and that the relation \((x, y) \equiv (x', y')\) is equivalent to the relation
\end{enumerate}

\[
f _ {a} (x) \left(f _ {a} \left(x ^ {\prime}\right)\right) ^ {- 1} = f _ {a} (y) \left(f _ {a} \left(y ^ {\prime}\right)\right) ^ {- 1}.
\]

For G to be commutative, it is necessary and sufficient that the relation \((x,y)\equiv (x',y')\) imply \((x,x')\equiv (y,y')\) .

¶ 25. Let E be a set and f a mapping of \(E^{m}\) into E; we write

\[
f (x _ {1}, x _ {2}, \dots , x _ {m}) = x _ {1} x _ {2} \dots x _ {m};
\]

suppose that \(f\) satisfies the following conditions:

\[
(x _ {1} x _ {2} \dots x _ {m}) x _ {m + 1} \dots x _ {2 m - 1} = x _ {1} (x _ {2} x _ {3} \dots x _ {m + 1}) x _ {m + 2} \dots x _ {2 m - 1}\tag{1}
\]

identically;

\begin{enumerate}
\def\labelenumi{(\arabic{enumi})}
\setcounter{enumi}{1}
\tightlist
\item
  For all \(a_1, a_2, \ldots, a_{m-1}\) , the mappings
\end{enumerate}

\[
x \mapsto x a _ {1} a _ {2} \dots a _ {m - 1}
\]

\[
x \mapsto a _ {1} a _ {2} \dots a _ {m - 1} x
\]

are bijections of E onto E.

\begin{enumerate}
\def\labelenumi{(\alph{enumi})}
\tightlist
\item
  Show that identically
\end{enumerate}

\[
(x _ {1} x _ {2} \dots x _ {m}) x _ {m + 1} \dots x _ {2 m - 1} = x _ {1} x _ {2} \dots x _ {i} (x _ {i + 1} \dots x _ {i + m}) x _ {i + m + 1} \dots x _ {2 m - 1}
\]

for every index \(i\) such that \(1 \leqslant i \leqslant m - 1\) (argue by induction on \(i\) , by considering the element

\[
\left(\left(x _ {1} x _ {2} \dots x _ {m}\right) x _ {m + 1} \dots x _ {2 m - 1}\right) a _ {1} a _ {2} \dots a _ {m - 1}).
\]

\begin{enumerate}
\def\labelenumi{(\alph{enumi})}
\setcounter{enumi}{1}
\tightlist
\item
  For all \(a_1, a_2, \ldots, a_{m-2}\) , there exists \(u \in E\) such that for all \(x\)
\end{enumerate}

\[
x = x a _ {1} a _ {2} \dots a _ {m - 2} u = u a _ {1} a _ {2} \dots a _ {m - 2} x.
\]

\begin{enumerate}
\def\labelenumi{(\alph{enumi})}
\setcounter{enumi}{2}
\tightlist
\item
  In the set \(\mathbf{E}^k\) of sequences \((u_1, u_2, \ldots, u_k)\) of \(k\) elements of \(\mathbf{E}\) ( \(1 \leqslant k \leqslant m - 1\) ) consider the equivalence relation \(\mathbf{R}_k\) which states: for all \(x_1, x_2, \ldots, x_{m-k}, u_1 u_2 \ldots u_k x_1 x_2 \ldots x_{m-k} = v_1 v_2 \ldots v_k x_1 x_2 \ldots x_{m-k}\) ; let \(\mathbf{E}_k\) denote the quotient set \(\mathbf{E}^k / \mathbf{R}_k\) and \(\mathbf{G}\) the sum set (Set Theory, II, §4, no. 8) of \(\mathbf{E}_1 = \mathbf{E}, \mathbf{E}_2, \ldots, \mathbf{E}_{m-1}\) . Let \(\alpha \in \mathbf{E}_i\) , \(\beta \in \mathbf{E}_j\) ; if \((u_1, u_2, \ldots, u_i)\) is a sequence of class \(\alpha\) , \((v_1, \ldots, v_j)\) a sequence of class \(\beta\) , consider the sequence
\end{enumerate}

\[
\left(u _ {1}, u _ {2}, \dots , u _ {i}, v _ {1}, \dots , v _ {j}\right)
\]

in \(\mathbf{E}^{i + j}\) if \(i + j < m\) , the sequence

\[
\left(u _ {1}, u _ {2}, \dots , u _ {i + j - m}, \left(u _ {i + j - m + 1} \dots u _ {i} v _ {1} v _ {2} \dots v _ {j}\right)\right)
\]

in \(E^{i+j-m+1}\) if \(i+j\geqslant m\) ; show that the class of this sequence in \(E_{i+j}\) (resp. \(E_{i+j-m+1}\) ) depends only on the classes \(\alpha\) and \(\beta\) ; it is denoted by \(\alpha.\beta\) . Show that a group law is thus defined on G, that \(H=E_{m-1}\) is a normal subgroup of G and that the quotient group G/H is cyclic of order m-1; prove finally that E is identical with a class (mod. H) which generates G/H and that \(x_{1}x_{2}\ldots x_{m}\) is just the composition, in the group G, of the sequence \((x_{1},x_{2},\ldots,x_{m})\) .

¶ 26. Let G, G' be two groups, \(f: G \to G'\) a mapping such that, for two arbitrary elements \(x, y\) of G, \(f(xy) = f(x)f(y)\) or \(f(xy) = f(y)f(x)\) . It is proposed to prove that \(f\) is a homomorphism of G into G' or a homomorphism of G into the opposite group G'0 (in other words, either \(f(xy) = f(x)f(y)\) for every ordered pair \((x, y)\) or \(f(xy) = f(y)f(x)\) for every ordered pair \((x, y)\) ).

\begin{enumerate}
\def\labelenumi{(\alph{enumi})}
\item
  Show that the set \(N = f(e')\) ( \(e'\) the identity element of \(G'\) ) is a normal subgroup of G and that f factors into \(G \xrightarrow{p} G/N \xrightarrow{g} G'\) , where g is injective. Attention may therefore be confined to the case where f is injective, which will be assumed throughout what follows.
\item
  Show that if \(xy = yx\) then \(f(x)f(y) = f(y)f(x)\) (consider \(f(x^2y)\) and \(f(x^2y^2)\) , expressing them in several ways).
\item
  Show that if \(f(xy) = f(x)f(y)\) then \(f(yx) = f(y)f(x)\) . (Show with the aid of (b) that attention may be confined to the case where \(xy \neq yx\) .)
\item
  Show that \(f(xyx) = \hat{f}(x)f(y)f(x)\) . (Use (a) or (b) according to whether \(xy = yx\) or \(xy \neq yx\) ; in the second case, consider \(f(x^2y)\) and \(f(yx^2)\) .)
\item
  Let A be the set of \(x \in G\) such that \(f(xy) = f(x)f(y)\) for all \(y \in G\) and B the set of \(x \in G\) such that \(f(xy) = f(y)f(x)\) for all \(y \in G\) . Show that \(A \neq G\) , \(B \neq G\) and \(A \cup B = G\) is an impossible situation. (By virtue of (b) there would then exist \(x, y\) in A such that \(f(xy) = f(x)f(y) \neq f(y)f(x)\) and \(u, v\) in B such that \(f(uv) = f(v)f(u) \neq f(u)f(v)\) . Deduce from this a contradiction by considering successively the two possibilities \(xu \in A\) , \(xu \in B\) ; in the first case consider \(f(xuv)\) and in the second \(f(yxu)\) .)
\end{enumerate}

The rest of the proof consists of proving that \(\mathbf{A} \cup \mathbf{B} \neq \mathbf{G}\) is impossible, arguing by reductio ad absurdum. If \(\mathbf{A} \cup \mathbf{B} \neq \mathbf{G}\) , there exist \(a, b, c\) in \(\mathbf{G}\) such that

\[
f (a b) = f (a) f (b) \neq f (b) f (a) \quad \text { and } \quad f (a c) = f (c) f (a) \neq f (a) f (c).
\]

\begin{enumerate}
\def\labelenumi{(\alph{enumi})}
\setcounter{enumi}{5}
\item
  Show that \(f(c)f(a)f(b) = f(b)f(a)f(c)\) . (Consider two cases according to whether \(bc \neq cb\) or \(bc = cb\) . In the first case, consider \(f(bac)\) ; in the second, use (b) and (c) and consider \(f(abc), f(bca)\) and \(f(bac)\) .)
\item
  Consider \(f(abac)\) and obtain a contradiction.
\end{enumerate}

\BourbakiExercisesSection

\begin{enumerate}
\def\labelenumi{\arabic{enumi}.}
\tightlist
\item
  In a finite group G, show that the number of conjugates (no. 4) of an element \(a \in G\) is equal to the index of the normalizer of a and is therefore a divisor of the order of G.
\end{enumerate}

¶ 2. *If G is a finite group of order n, the number of automorphisms of G is \(\leqslant n^{\log n/\log 2}\) (show that there exists a generating set \(\{a_{1}, a_{2}, \ldots, a_{m}\}\) of G such that \(a_{i}\) does not belong to the subgroup generated by \(a_{1}, a_{2}, \ldots, a_{i-1}\) for \(2 \leqslant i \leqslant m\) ; deduce that \(2^{m} \leqslant n\) and that the number of automorphisms of G is \(\leqslant n^{m})\cdot*\)

\begin{enumerate}
\def\labelenumi{\arabic{enumi}.}
\setcounter{enumi}{2}
\tightlist
\item
  Let \(\Gamma\) be the group of automorphisms of a group \(G\) and \(\Delta\) the group of inner automorphisms of \(G\) ; show that \(\Delta\) is a normal subgroup of \(\Gamma\) . For an automorphism \(\sigma\) of \(G\) to be permutable with all the inner automorphisms of \(G\) , it is necessary and sufficient that, for all \(x \in G\) , \(x^{-1}\sigma(x)\) belong to the centre of \(G\) ; deduce that, if the centre of \(G\) is reduced to the identity element, so is the centralizer of \(\Delta\) in \(\Gamma\) .
\end{enumerate}

¶ 4. (a) Let G be a non-commutative simple group, Γ the automorphism group of G and Δ the group of inner automorphisms of G (isomorphic to G). If s is an automorphism of the group Γ, show that \(s(\Delta) = \Delta\) (using Exercise 3 above and § 4, no. 9, Proposition 15, note that the intersection \(\Delta \cap s(\Delta)\) cannot consist only of the identity element of Γ).

\begin{enumerate}
\def\labelenumi{(\alph{enumi})}
\setcounter{enumi}{1}
\item
  Show that the only automorphism of \(\Gamma\) leaving invariant each of the elements of \(\Delta\) is the identity element (use the fact that this automorphism leaves invariant \(\alpha\) and \(\sigma \alpha \sigma^{-1}\) for all \(\alpha \in \Delta\) and \(\sigma \in \Gamma\) and use Exercise 3).
\item
  Let \(s\) be an automorphism of \(\Gamma\) , \(\phi\) the isomorphism \(x \mapsto \operatorname{Int}(x)\) of G onto \(\Delta\) , \(\psi\) the inverse isomorphism and \(\sigma\) the automorphism \(\psi \circ s \circ \phi\) of G; show that the automorphism \(\xi \mapsto \sigma^{-1}s(\xi)\sigma\) of \(\Gamma\) is the identity automorphism (use (b), noting that, for all \(x \in G\) , \(s(\operatorname{Int}(x)) = \operatorname{Int}(\sigma(x)))\) . Deduce that every automorphism of \(\Gamma\) is an inner automorphism.
\end{enumerate}

5. Let G be a group.

\begin{enumerate}
\def\labelenumi{(\alph{enumi})}
\item
  If H is a subgroup of G of finite index n, show that the intersection N of the conjugates of H is of index a divisor of \(n!\) (note that G/N is isomorphic to a subgroup of \(S_{n}\) ).
\item
  G is called residually finite, if the intersection of its subgroups of finite index is \(\{e\}\) . Show that this is equivalent to saying that G is isomorphic to a subgroup of a product of finite groups.
\end{enumerate}
